% !TeX root = ../../Book1.tex
% 纲要标题：## 附录 D　G-结构与 Cartan 几何
% 纲要位置：几何学计划纲要.md，第 981 行。
% 正文状态：初稿完成，含完整证明、模型计算及习题解答；尚未正式审核。

% 纲要内容开始
% 
% > 沿标架丛约化发展结构方程；选定低阶模型给出完整计算，一般等价算法另属专题。
% 
% * **结构群约化与内蕴挠率**
%   * G-结构及适应标架
%   * 联络选择与内蕴挠率
%   * 正交、几乎复及接触模型
% 
% * **Cartan 联络与齐性模型**
%   * Cartan 联络的公理
%   * 曲率、平坦性及路径展开
%   * Cartan 联络与主联络
% 
% * **标架归一化与几何等价**
%   * 结构方程中的规范自由
%   * 曲率梯度与曲面标架的归一化
%   * 局部等价与整体单值
% 
% ---
% 
% 纲要内容结束

\chapter{G-结构与 Cartan 几何}
\label{g1:app:D}
\BookChapterNumber{g1:app:D}

\begin{chapterabstract}
	几何结构可以通过允许哪些标架来记录，而这些标架随位置变化时，又会产生挠率与曲率。这个附录把结构群约化推进到微分层面：先用联络的差说明内蕴挠率为何与联络选择无关，再通过正交、几乎复及接触模型比较不同约化的障碍。随后定义以齐性空间为模型的 Cartan 联络，证明其平坦性与局部模型实现的关系，并解释它和主丛水平联络的区别。最后以曲面曲率梯度固定标架，计算剩余结构函数，说明局部等价的检验及整体单值所记录的信息。
\end{chapterabstract}

\begin{learninggoals}
	\begin{itemize}
		\item 用切丛标架的约化描述几何结构，写出适应标架上的典范形式及其变换规律。
		\item 计算改变相容联络时挠率的变化，并在商空间中识别不能消去的部分。
		\item 比较正交、几乎复和接触模型，区分无挠、平坦与具有局部标准形。
		\item 核对 Cartan 联络的三项公理，计算其曲率，并说明平坦模型的局部构造。
		\item 通过展开方程比较路径与齐性模型，辨别 Cartan 联络和 Ehresmann 主联络所提供的数据。
		\item 在曲率梯度非零处归一化曲面标架，用结构方程检验局部等价，并辨认整体周期障碍。
	\end{itemize}
\end{learninggoals}

在平面开集上取两条处处独立的一形式
\[
	\theta^1=\mathrm{d}x,\qquad \theta^2=e^{x^2}\,\mathrm{d}y.
\]
它们构成一个余标架，因而可以用来记录切向量的两个分量。但若希望它们分别是两个新坐标的微分，就遇到了
\[
	\mathrm{d}\theta^2=2x\,\theta^1\wedge\theta^2.
\]
在 \(x\ne0\) 处，第二条形式不是闭形式，不能写成局部坐标函数的微分。把它乘以 \(e^{-x^2}\) 虽然恢复了 \(\mathrm{d}y\)，却改变了度量 \(g=(\theta^1)^2+(\theta^2)^2\)。如果希望保留这个度量，就只能用正交矩阵改变正交余标架，不能任意伸缩其中一条形式。

\ProseParagraphBreak

这说明“选择方便的标架”有一个先于计算的问题：哪些变化仍然保留所研究的结构？\chapref{g1:ch:12}和\chapref{g1:ch:18}用结构群约化回答了逐点的选择问题。现在还要比较这些标架在邻近点的变化。相容联络提供一种比较方法，但联络往往不唯一；改变联络能够消去哪些挠率分量，又有哪些分量属于几何结构本身，是本附录的第一个问题。

\ProseParagraphBreak

即使可以把挠率全部消去，几何也未必成为Euclid 几何。曲面上的无挠度量联络仍可有曲率。若把“允许的转动”和“沿曲面前进的位移”放入同一个Lie 代数值形式，便得到一种同时记录标架及切向运动的描述。Cartan 联络把这种构造推广到一般齐性模型；其平坦性表达局部与模型一致，整体是否能用同一份模型坐标描述，还取决于绕行和延拓。

\ProseParagraphBreak

本附录使用\chapref{g1:ch:12}的标架丛、关联丛与光滑单位分解，\chapref{g1:ch:13}的Lie 括号和 Frobenius 定理，以及\chapref{g1:ch:15}的挠率、曲率与结构方程。主丛、局部截面、联络形式和群值微分方程采用\chapref{g1:ch:18}的右作用约定；齐性空间的光滑结构接续\chapref{g1:ch:17}。曲面模型只需恢复\chapref{g1:ch:09}的正向正交标架及 \(\mathrm{d}\alpha=-K\,\mathrm{d}A\) 约定。各节将从所用对象的定义出发，把结构等价的问题化成适应标架与结构函数的比较。

\section{结构群约化与内蕴挠率}
\label{g1:sec:D01}

第十二章把度量、几乎复结构等逐点数据组织成标架丛的结构群约化，第十八章又说明主联络怎样在不同纤维之间作比较。把两件事放在一起，会遇到一个新的问题：既要求联络保持指定结构，又要求它无挠，这两项能否同时做到？第十五章对任意切丛联络作修正可以消去挠率，但修正后的联络未必还保持原来的几何结构。本节把允许的修正写成一个有限维线性映射，再在目标空间中对该映射的像取商，提取不随联络改变的部分。

\subsection{G-结构及适应标架}
\label{g1:subsec:D01:01}

本节中，\(M\) 是有限维、无边界、Hausdorff、第二可数的光滑流形，\(\dim M=n\geq1\)，\(V=\mathbb R^n\)，\(G\subset\operatorname{GL}(V)\) 是闭Lie 子群，Lie 代数为 \(\mathfrak g\)。沿用\cref{g1:def:frame-bundle}的约定，\(x\) 处的标架是线性同构 \(u:V\to\mathrm{T}_xM\)，右作用为 \(u g=u\circ g\)。指定标准基以后，\(u\) 也就是一个有序切向量基。

\begin{definition}
	\label{g1:def:g-structure}
	设 \(n\ge1\)，\(M\) 为 \(n\) 维、无边界、Hausdorff、第二可数光滑流形，\(V=\mathbb R^n\)，\(G\subset\operatorname{GL}(V)\) 为闭 Lie 子群。标架 \(u:V\to\mathrm T_xM\) 的右作用为 \(ug=u\circ g\)。\(M\) 上的一个 \term[Gjiegou]{\(G\)-结构}[\(G\)-structure]是标架丛中的光滑主 \(G\) 子丛
	\[
		\pi:P\subset\operatorname{Fr}(\mathrm{T}M)\longrightarrow M.
	\]
	它的局部截面称为\term[shiyingbiaojia]{适应标架}[adapted frame]。具体说，存在开覆盖 \(\{U_i\}\) 及光滑标架 \(s_i:U_i\to P\)，使
	\[
		P|_{U_i}=\{s_i(x)g:x\in U_i,\ g\in G\},
		\qquad s_j=s_i g_{ij},\quad g_{ij}:U_i\cap U_j\to G.
	\]
	这里要求局部表达 \((x,g)\mapsto s_i(x)g\) 给出主丛平凡化。
\end{definition}

这正是\cref{g1:def:frame-structure-reduction}在切丛上的主丛表述。若从局部适应标架及其 \(G\) 值过渡函数出发，上述集合在交叠处一致；用 \(U_i\times G\) 的坐标粘合，就得到所需主子丛。反过来，主子丛的局部截面给出这样的标架。因此，局部适应标架及其光滑的过渡关系共同确定一个约化。

\ProseParagraphBreak

标架还提供了仅由主丛结构无法得到的信息：它把固定空间 \(V\) 与底空间的切空间联系起来。这个联系可以记在 \(P\) 上的一形式中。对 \(u\in P_x\)、\(\xi\in\mathrm{T}_uP\)，定义
\begin{equation}
	\label{g1:eq:tautological-form}
	\theta_u(\xi)=u^{-1}\bigl(\mathrm{d}\pi_u\xi\bigr)\in V.
\end{equation}
称 \(\theta\in\Omega^1(P;V)\) 为标架约化上的\term[dianfanyixingshi]{典范一形式}[tautological form]。\symidx{\(\theta\)：标架丛及其约化上的典范一形式}它先将总空间切向量投到底空间，再用当前标架读取分量；这个定义没有选择联络。

\begin{proposition}[典范形式与余标架]
	\label{g1:prop:tautological-form-properties}
	设 \(M\) 为 \(n\) 维光滑流形，\(V=\mathbb R^n\)，\(G\subset\operatorname{GL}(V)\) 为闭 Lie 子群，\(\pi:P\subset\operatorname{Fr}(\mathrm{T}M)\to M\) 为 \(G\)-结构。对标架 \(u:V\to\mathrm T_xM\)，定义典范形式 \(\theta_u(\xi)=u^{-1}(\mathrm d\pi_u\xi)\)，并以 \(R_g(u)=ug\) 记右作用。则 \(\theta\) 光滑，每个 \(\theta_u:\mathrm{T}_uP\to V\) 满射，且
	\[
		\ker\theta_u=\ker\mathrm{d}\pi_u,\qquad
		R_g^*\theta=g^{-1}\theta\quad(g\in G).
	\]
	若 \(s=(e_1,\ldots,e_n)\) 是局部适应标架，则 \(s^*\theta\) 是其对偶余标架组成的 \(V\) 值一形式。
\end{proposition}
\begin{proof}
	在局部表达 \(u=s(x)g\) 下，\eqref{g1:eq:tautological-form}成为 \(g^{-1}s(x)^{-1}\mathrm{d}\pi\)，其矩阵系数光滑，所以 \(\theta\) 光滑。\(\pi\) 为淹没且 \(u^{-1}\) 可逆，故 \(\theta_u\) 满射，核恰为垂直空间。又因 \(\pi\circ R_g=\pi\)，有
	\[
		\theta_{ug}(\mathrm{d}R_g\xi)
		=(ug)^{-1}\mathrm{d}\pi_u\xi
		=g^{-1}\theta_u(\xi).
	\]
	最后，\(\pi\circ s=\operatorname{id}\) 给出 \((s^*\theta)_x(X)=s(x)^{-1}X\)。把 \(X=\sum_iX^ie_i\) 代入，结果正是列向量 \((X^1,\ldots,X^n)^{\mathsf T}\)，所以各分量为对偶余标架。
\end{proof}

在底空间上，\(s^*\theta\) 的 \(n\) 个分量确是余标架；在 \(P\) 上，\(\theta\) 却有垂直核，通常不是总空间的余标架，因为 \(\dim P=n+\dim G\)。它也不是主联络形式：典范形式取值于 \(V\)，并在所有垂直向量上为零；主联络形式取值于 \(\mathfrak g\)，需要读出垂直方向中的群运动。下一节的 Cartan 联络将同时处理这两种信息。

\subsection{联络选择与内蕴挠率}
\label{g1:subsec:D01:02}

\cref{g1:thm:principal-connection-existence}保证 \(P\) 上存在主联络。固定一个主联络形式 \(\omega\in\Omega^1(P;\mathfrak g)\)，其公理为
\[
	\omega(A^\#)=A,\qquad R_g^*\omega=\operatorname{Ad}_{g^{-1}}\omega.
\]
这里 \(A^\#_u=\left.\frac{\mathrm{d}}{\mathrm{d}t}\right|_{0}u\exp(tA)\) 是 \(A\in\mathfrak g\) 产生的垂直基本向量。标准表示给出
\[
	P\times_G V\longrightarrow\mathrm{T}M,\qquad [u,v]\longmapsto u(v).
\]
它与\cref{g1:prop:frame-associated-equivalence}中的映射相同：\([ug,g^{-1}v]\) 的像仍为 \(u(v)\)，在适应标架中则是逐纤维恒等坐标，故为光滑向量丛同构。由\cref{g1:thm:associated-covariant-derivative}，\(\omega\) 因而诱导切丛联络。若 \(a=s^*\omega\)，则其局部公式为 \(\nabla(s f)=s(\mathrm{d}f+a f)\)。这样的联络称为适应该结构的 \(G\)-联络；局部联络矩阵取值于 \(\mathfrak g\)，正是允许的标架变换所要求的限制。

\begin{definition}
	\label{g1:def:g-structure-torsion-form}
	设 \(M\) 为 \(n\) 维光滑流形，\(V=\mathbb R^n\)，\(G\subset\operatorname{GL}(V)\) 为闭 Lie 子群，Lie 代数为 \(\mathfrak g\)，\(\pi:P\subset\operatorname{Fr}(\mathrm TM)\to M\) 为 \(G\)-结构。令 \(\omega\in\Omega^1(P;\mathfrak g)\) 为主 \(G\) 联络，\(\theta_u(\xi)=u^{-1}(\mathrm d\pi_u\xi)\) 为典范 \(V\) 值形式。定义 \(\omega\) 的挠率形式为
	\begin{equation}
		\label{g1:eq:g-structure-torsion-form}
		T^\omega=\mathrm{d}\theta+\omega\wedge\theta\in\Omega^2(P;V),
	\end{equation}
	其中矩阵对向量的作用约定为
	\[
		(\omega\wedge\theta)(\xi,\eta)
		=\omega(\xi)\theta(\eta)-\omega(\eta)\theta(\xi).
	\]
\end{definition}

\begin{proposition}[挠率的水平性]
	\label{g1:prop:g-torsion-horizontal}
	设 \(M\) 为 \(n\) 维光滑流形，\(V=\mathbb R^n\)，\(G\subset\operatorname{GL}(V)\) 为闭 Lie 子群，\(\pi:P\subset\operatorname{Fr}(\mathrm TM)\to M\) 为 \(G\)-结构，\(\omega\) 为其主 \(G\) 联络。令 \(\theta_u=u^{-1}\mathrm d\pi_u\)、\(T^\omega=\mathrm d\theta+\omega\wedge\theta\)，并通过 \([u,v]\mapsto u(v)\) 将标准表示诱导的关联丛联络视为切丛联络 \(\nabla\)。则 \(T^\omega\) 是水平且等变的，即只要一个输入为垂直向量，其值便为零，并且 \(R_g^*T^\omega=g^{-1}T^\omega\)。对光滑向量场 \(X,Y\)，它对应的切丛值二形式正是诱导联络的挠率
	\[
		T^\nabla(X,Y)=\nabla_XY-\nabla_YX-[X,Y].
	\]
\end{proposition}
\begin{proof}
	由典范形式的等变性，沿 \(A^\#\) 的流求导得到 \(\mathcal L_{A^\#}\theta=-A\theta\)。又有 \(\iota_{A^\#}\theta=0\)，所以对每个分量使用 Cartan 公式，得
	\[
		\iota_{A^\#}\mathrm{d}\theta=-A\theta,
		\qquad
		\iota_{A^\#}(\omega\wedge\theta)=A\theta.
	\]
	两项相消，而各点的垂直空间由基本向量填满，故 \(T^\omega\) 水平。对于固定 \(g\)，外微分与拉回交换，且
	\[
		(g^{-1}\omega g)\wedge(g^{-1}\theta)
		=g^{-1}(\omega\wedge\theta),
	\]
	所以得出等变性。

	再令 \(\vartheta=s^*\theta\)、\(a=s^*\omega\)。拉回后
	\[
		s^*T^\omega=\mathrm{d}\vartheta+a\wedge\vartheta.
	\]
	对底空间向量场 \(X,Y\)，其右边等于
	\[
		X(\vartheta(Y))-Y(\vartheta(X))-\vartheta([X,Y])
		+a(X)\vartheta(Y)-a(Y)\vartheta(X).
	\]
	利用局部协变导数公式，这就是 \(\vartheta(\nabla_XY-\nabla_YX-[X,Y])\)。它同时复现了\cref{g1:prop:first-structure-equation}中的第一结构方程。
\end{proof}

水平性使挠率可以完全由底空间方向的标架分量描述。令
\[
	W=\Lambda^2V^*\otimes V.
\]
每个 \(u\in P\) 唯一确定 \(t^\omega(u)\in W\)，使
\begin{equation}
	\label{g1:eq:g-torsion-function}
	T^\omega_u(\xi,\eta)
	=t^\omega(u)\bigl(\theta_u\xi,\theta_u\eta\bigr).
\end{equation}
确实，\(\theta_u\) 满射；对给定 \(v,w\in V\) 任选提升，两个不同提升之差垂直，故 \(T^\omega\) 的值相同。在局部适应标架下，此定义就是把光滑挠率张量取分量，因而 \(t^\omega\) 光滑。称它为相对于联络 \(\omega\) 的挠率函数。

\ProseParagraphBreak

现在改变联络。由\cref{g1:prop:principal-connection-affine}，两个主联络之差 \(\beta=\omega'-\omega\) 是水平的伴随等变一形式。因此唯一存在光滑函数
\[
	A:P\longrightarrow C\coloneqq V^*\otimes\mathfrak g,
	\qquad \beta_u=A(u)\circ\theta_u.
\]
这里 \(A(u)\) 把一个底空间方向的标架分量映成允许的无穷小标架变换。把这个变换依次作用在另一个方向上，再作交替化，就得到线性映射
\begin{equation}
	\label{g1:eq:intrinsic-torsion-delta}
	\delta:C\longrightarrow W,\qquad
	(\delta A)(v,w)=A(v)w-A(w)v.
\end{equation}

\begin{proposition}[改变联络时的挠率变化]
	\label{g1:prop:g-torsion-change}
	设 \(M\) 为 \(n\) 维光滑流形，\(V=\mathbb R^n\)，\(G\subset\operatorname{GL}(V)\) 为闭 Lie 子群，Lie 代数为 \(\mathfrak g\)，\(\pi:P\subset\operatorname{Fr}(\mathrm TM)\to M\) 为 \(G\)-结构，\(\theta_u=u^{-1}\mathrm d\pi_u\) 为典范形式。对 \(P\) 上的两主联络 \(\omega,\omega'\)，设 \(A:P\to V^*\otimes\mathfrak g\) 满足 \((\omega'-\omega)_u=A(u)\circ\theta_u\)。对任一主联络 \(\eta\)，令 \(T^\eta=\mathrm d\theta+\eta\wedge\theta\)，并由
	\[
		T^\eta_u(\xi,\zeta)=t^\eta(u)(\theta_u\xi,\theta_u\zeta)
	\]
	定义 \(t^\eta:P\to\Lambda^2V^*\otimes V\)。定义交替化映射 \((\delta B)(v,w)=B(v)w-B(w)v\)，其中 \(B\in V^*\otimes\mathfrak g\)。则
	\begin{equation}
		\label{g1:eq:g-torsion-change}
		t^{\omega'}=t^\omega+\delta A.
	\end{equation}
	在表示
	\[
		(g\cdot t)(v,w)=g\,t(g^{-1}v,g^{-1}w),
		\qquad
		(g\cdot A)(v)=g A(g^{-1}v)g^{-1}
	\]
	下，\(\delta\) 是 \(G\) 等变映射，而 \(t^\omega(ug)=g^{-1}\cdot t^\omega(u)\)、\(A(ug)=g^{-1}\cdot A(u)\)。
\end{proposition}
\begin{proof}
	将两个挠率公式相减，得 \(T^{\omega'}-T^\omega=\beta\wedge\theta\)。在 \(\theta\xi=v\)、\(\theta\eta=w\) 的两个输入上，右边为 \(A(v)w-A(w)v\)，得到\eqref{g1:eq:g-torsion-change}。

	等变性可直接核对。例如
	\[
		\bigl(\delta(g\cdot A)\bigr)(v,w)
		=g\bigl(A(g^{-1}v)g^{-1}w-A(g^{-1}w)g^{-1}v\bigr)
		=(g\cdot\delta A)(v,w).
	\]
	\(T^\omega\) 与 \(\theta\) 的等变性给出
	\[
		t^\omega(ug)(v,w)=g^{-1}t^\omega(u)(gv,gw),
	\]
	即所述公式。对 \(\beta=A\theta\) 使用 \(R_g^*\beta=g^{-1}\beta g\) 同样得到 \(A\) 的变换式。
\end{proof}

\begin{definition}
	\label{g1:def:intrinsic-torsion}
	设 \(n\ge1\)，\(M\) 为 \(n\) 维、无边界、Hausdorff、第二可数光滑流形，\(V=\mathbb R^n\)，\(G\subset\operatorname{GL}(V)\) 为闭 Lie 子群，Lie 代数为 \(\mathfrak g\)，\(\pi:P\subset\operatorname{Fr}(\mathrm TM)\to M\) 为 \(G\)-结构。任取主 \(G\) 联络 \(\omega\)，以 \(\theta_u=u^{-1}\mathrm d\pi_u\) 记典范形式，并由
	\[
		(\mathrm d\theta+\omega\wedge\theta)_u(\xi,\zeta)
		=t^\omega(u)(\theta_u\xi,\theta_u\zeta)
	\]
	定义挠率函数 \(t^\omega:P\to\Lambda^2V^*\otimes V\)。令 \(\delta:V^*\otimes\mathfrak g\to\Lambda^2V^*\otimes V\) 满足 \((\delta A)(v,w)=A(v)w-A(w)v\)，并记
	\[
		W_G=\frac{\Lambda^2V^*\otimes V}{\delta(V^*\otimes\mathfrak g)}.
	\]
	\(G\)-结构 \(P\) 的\term[neiyunnaolv]{内蕴挠率}[intrinsic torsion]是等变函数
	\[
		\tau_P:P\longrightarrow W_G,\qquad
		\tau_P(u)=[t^\omega(u)],
	\]
	此商类与主联络 \(\omega\) 的选择无关，也可视为关联丛 \(P\times_G W_G\to M\) 的光滑截面。\symidx{\(\tau_P\)：\(G\)-结构的内蕴挠率}
\end{definition}

由\cref{g1:prop:g-torsion-change}，更换联络只会增加 \(\delta\) 的像中的项，因此这个商类与联络无关。\(\delta\) 等变又保证其像为不变子空间，商空间具有良定义的 \(G\) 作用；换适应标架时，各局部商类正按这个表示变换。因此内蕴挠率是该约化本身的几何数据。它是商空间中的类；若要用 \(W\) 中某个补空间的张量来代表它，需再选取与换标架相容的补空间。

\begin{theorem}[无挠适应联络的存在]
	\label{g1:thm:intrinsic-torsion-zero}
	设 \(n\ge1\)，\(M\) 为 \(n\) 维、无边界、Hausdorff、第二可数光滑流形，\(G\subset\operatorname{GL}(n,\mathbb R)\) 为闭Lie 子群，\(\pi:P\subset\operatorname{Fr}(\mathrm{T}M)\to M\) 为 \(G\)-结构。记 \(V=\mathbb R^n\)、\(\mathfrak g=\operatorname{Lie}(G)\)，令 \(\delta:V^*\otimes\mathfrak g\to\Lambda^2V^*\otimes V\) 满足 \((\delta A)(v,w)=A(v)w-A(w)v\)。对任一主联络 \(\omega\)，以 \(\theta_u=u^{-1}\mathrm d\pi_u\) 记典范形式，由 \((\mathrm d\theta+\omega\wedge\theta)_u(\xi,\zeta)=t^\omega(u)(\theta_u\xi,\theta_u\zeta)\) 定义挠率函数，并令 \(\tau_P(u)=[t^\omega(u)]\in\operatorname{coker}\delta\)。则 \(\tau_P=0\) 当且仅当 \(P\) 存在无挠 \(G\)-联络。若存在，则所有这样的联络组成仿射空间，其差由 \(P\times_G\ker\delta\) 的光滑截面给出。
\end{theorem}
\begin{proof}
	若存在无挠联络，其挠率函数为零，故内蕴挠率为零。反向，先由\cref{g1:thm:principal-connection-existence}取一个全局主联络 \(\omega\)。若 \(\tau_P=0\)，则 \(t^\omega(u)\in\operatorname{im}\delta\) 对每个 \(u\) 成立。

	取局部适应标架 \(s_i:U_i\to P\)。有限维满射 \(\delta:C\to\operatorname{im}\delta\) 有线性右逆 \(R:\operatorname{im}\delta\to C\)：在 \(C\) 中选取 \(\ker\delta\) 的一个线性补空间，将 \(\delta\) 在该补空间上的限制求逆即可。定义
	\[
		A_i(s_i(x))=-R\bigl(t^\omega(s_i(x))\bigr),
		\qquad
		A_i(s_i(x)g)=g^{-1}\cdot A_i(s_i(x)).
	\]
	平凡化中的 \((x,g)\) 唯一，所以后一式给出光滑、等变的延拓。由 \(\delta\) 等变性，\(\delta A_i=-t^\omega\) 在整个 \(P|_{U_i}\) 成立。于是 \(\beta_i=A_i\theta\) 水平且伴随等变，\(\omega_i=\omega+\beta_i\) 是 \(P|_{U_i}\) 上的主联络，且由\eqref{g1:eq:g-torsion-change}无挠。这里右逆 \(R\) 只用于一个局部平凡化，并没有要求它对 \(G\) 等变。

	取从属于 \(\{U_i\}\) 的局部有限光滑单位分解 \(\{\chi_i\}\)，其支集包含于相应 \(U_i\)。将 \((\chi_i\circ\pi)\beta_i\) 在 \(P|_{U_i}\) 之外补零；支集条件保证补零后仍光滑。定义
	\[
		\omega_0=\omega+\sum_i(\chi_i\circ\pi)\beta_i.
	\]
	该和局部有限，各项均水平且伴随等变，故 \(\omega_0\) 是全局主联络。在每一点附近也可写成仿射组合 \(\sum_i(\chi_i\circ\pi)\omega_i\)，权重和为一，因而
	\[
		T^{\omega_0}
		=\sum_i(\chi_i\circ\pi)
		\bigl(\mathrm{d}\theta+\omega_i\wedge\theta\bigr)=0.
	\]
	这一步没有权重的导数项，因为挠率公式中对 \(\theta\) 求外微分，对联络形式只作代数运算。于是得到全局无挠适应联络。

	最后，固定一个无挠联络后，另一个联络由 \(\beta=A\theta\) 给出。它也无挠恰好等价于 \(\delta A=0\)。等变函数 \(A:P\to\ker\delta\) 与关联丛 \(P\times_G\ker\delta\) 的截面一一对应，这证明最后一项。
\end{proof}

这个结论同时区分了存在性和唯一性：商空间中的类控制能否消去挠率，\(\ker\delta\) 则控制消去以后还有多少联络可选。即使无挠联络存在，它的曲率也未必为零；即使某类几何具有统一的局部标准模型，该模型的内蕴挠率也未必为零。下面的例子将分别说明这些区别。

\subsection{正交、几乎复及接触模型}
\label{g1:subsec:D01:03}

先设 \(g\) 是 \(M\) 上的Riemann 度量。由\cref{g1:prop:metric-orthogonal-reduction}，正交单位标架组成 \(\operatorname{O}(n)\)-结构。在这样的标架中，度量矩阵为 \(I\)，\cref{g1:eq:metric-connection-matrix}化为 \(a^{\mathsf T}+a=0\)。因此 \(\operatorname{O}(n)\)-联络恰为保持 \(g\) 的度量联络。

\begin{proposition}[正交约化的挠率修正]
	\label{g1:prop:orthogonal-delta-isomorphism}
	设 \(V\) 是 \(n\) 维实内积空间。线性映射
	\[
		\delta:V^*\otimes\mathfrak{so}(V)\longrightarrow\Lambda^2V^*\otimes V
	\]
	由 \((\delta A)(v,w)=A(v)w-A(w)v\) 定义，则它是同构。因此，任意 \(n\) 维光滑流形上的 Riemann 度量所确定的正交标架约化具有恒为零的内蕴挠率，并且有唯一的无挠适应联络。
\end{proposition}
\begin{proof}
	先证单射。设 \(\delta A=0\)，令
	\[
		a(v,w,z)=\langle A(v)w,z\rangle.
	\]
	\(\delta A=0\) 说明 \(a\) 的前两个位置对称；\(A(v)\in\mathfrak{so}(V)\) 说明后两个位置反对称。依次使用这两种关系，得到
	\[
	\begin{aligned}
		a(v,w,z)&=a(w,v,z)=-a(w,z,v)=-a(z,w,v)\\
		&=a(z,v,w)=a(v,z,w)=-a(v,w,z).
	\end{aligned}
	\]
	故 \(a=0\)，由内积非退化性得 \(A=0\)。两边的维数均为 \(n^2(n-1)/2\)，所以单射即为同构。于是 \(W_{\operatorname{O}(n)}=0\) 且 \(\ker\delta=0\)，由\cref{g1:thm:intrinsic-torsion-zero}分别得到存在性与唯一性。
\end{proof}

这也从结构群约化的角度解释了\chapref{g1:ch:19}所用的 Levi--Civita 联络：斜对称矩阵的自由度恰能唯一地修正全部挠率。内蕴挠率在所有Riemann 度量上都为零，所以它不负责区分平面、球面等不同曲率的几何；那些差别还须由这个唯一联络的曲率来测量。

\ProseParagraphBreak

设 \(J:\mathrm{T}M\to\mathrm{T}M\) 光滑且 \(J^2=-I\)，\(\dim M=2m\)。\cref{g1:prop:complex-frame-reduction}给出 \(\operatorname{GL}(m,\mathbb C)\) 约化，相应的实适应标架可以写成 \(e_1,Je_1,\ldots,e_m,Je_m\)。在适应标架中 \(J\) 是常矩阵 \(J_0\)，所以
\[
	\nabla J=0\quad\Longleftrightarrow\quad aJ_0=J_0a.
\]
右边恰好要求 \(a\) 取值于 \(\mathfrak{gl}(m,\mathbb C)\) 的实表示。因此这里的适应联络就是保持 \(J\) 的联络。

\ProseParagraphBreak

复化后的 Lie 括号直接给出一种在改变适应联络时保持不变的挠率分量。把切丛和 \(J\) 复化，令
\[
	\mathrm{T}^{1,0}M=\ker(J-\mathrm iI),\qquad
	\mathrm{T}^{0,1}M=\ker(J+\mathrm iI).
\]
投影 \(\operatorname{pr}^{1,0}=\tfrac12(I-\mathrm iJ)\)、\(\operatorname{pr}^{0,1}=\tfrac12(I+\mathrm iJ)\) 光滑且互补，两个像均为复秩 \(m\) 的子丛。对局部 \(\mathrm{T}^{1,0}M\) 截面 \(Z,W\)，定义
\begin{equation}
	\label{g1:eq:almost-complex-bracket-obstruction}
	\mathcal K_J(Z,W)=\operatorname{pr}^{0,1}[Z,W].
\end{equation}
这里括号按复双线性延拓。若在输入中乘以光滑复函数，括号的额外导数项仍在 \(\mathrm{T}^{1,0}M\)，投影后消失；例如 \(\operatorname{pr}^{0,1}[fZ,W]=f\operatorname{pr}^{0,1}[Z,W]\)。所以 \(\mathcal K_J\) 是良定义的光滑张量，测量两个 \((1,0)\) 方向的括号落出该子丛的部分。

\begin{proposition}[几乎复结构的括号障碍]
	\label{g1:prop:almost-complex-intrinsic-component}
	设 \(M\) 为实 \(2m\) 维光滑流形，\(J\in\Gamma(\operatorname{End}(\mathrm TM))\) 光滑且 \(J^2=-I\)，\(\nabla\) 为 \(\mathrm TM\) 上满足 \(\nabla J=0\) 的线性联络，挠率为 \(T^\nabla(X,Y)=\nabla_XY-\nabla_YX-[X,Y]\)。复化切丛后，令 \(\mathrm T^{1,0}M=\ker(J-\mathrm iI)\)、\(\operatorname{pr}^{0,1}=(I+\mathrm iJ)/2\)，并对光滑截面 \(Z,W\) 定义 \(\mathcal K_J(Z,W)=\operatorname{pr}^{0,1}[Z,W]\)。将联络及其挠率复化，则对 \(Z,W\in\Gamma(\mathrm{T}^{1,0}M)\)，有
	\[
		\operatorname{pr}^{0,1}T^\nabla(Z,W)=-\mathcal K_J(Z,W).
	\]
	因此右边是其 \(\operatorname{GL}(m,\mathbb C)\)-结构的内蕴挠率中不随适应联络改变的一个分量；若 \(\mathcal K_J\neq0\)，则不存在无挠适应联络。
\end{proposition}
\begin{proof}
	因为 \(\nabla J=0\)，由 \(JW=\mathrm iW\) 求协变导数可得 \(J\nabla_ZW=\mathrm i\nabla_ZW\)。所以 \(\nabla_ZW\) 与 \(\nabla_WZ\) 都在 \(\mathrm{T}^{1,0}M\) 中。将挠率恒等式投影到 \(\mathrm{T}^{0,1}M\)，这两项消失，只留下 \(-\operatorname{pr}^{0,1}[Z,W]\)。这也说明改变适应联络所产生的 \(\delta A\) 在这个投影下为零，故该分量确由内蕴商类决定。
\end{proof}

若 \(J\) 来自一组复坐标，那么 \(\mathrm{T}^{1,0}M\) 局部由相互对易的复坐标向量场张成，其括号仍落在本子丛，故 \(\mathcal K_J=0\)。反过来从括号条件构造复坐标，需要额外的可积性定理；非零的括号障碍则已经排除了复坐标的存在，下面就计算一个这样的例子。

\begin{example}[不能消去的几乎复挠率分量]
	\label{g1:exm:almost-complex-nonzero-intrinsic}
	在 \(\mathbb R^4\) 的坐标 \((x,y,z,t)\) 中取全局标架
	\[
		e_1=\partial_x,\quad e_2=\partial_y,\quad
		e_3=\partial_z,\quad e_4=\partial_t+x\partial_y.
	\]
	规定 \(Je_1=e_2\)、\(Je_2=-e_1\)、\(Je_3=e_4\)、\(Je_4=-e_3\)，得到光滑几乎复结构。唯一非零的基本括号为 \([e_1,e_4]=e_2\) 及其反向。令
	\[
		Z=e_1-\mathrm i e_2,\qquad W=e_3-\mathrm i e_4,
	\]
	它们属于 \(\mathrm{T}^{1,0}M\)，且
	\[
		[Z,W]=-\mathrm i e_2=\frac12(Z-\overline Z),
		\qquad \mathcal K_J(Z,W)=-\frac12\overline Z\neq0.
	\]
	因此此结构没有无挠适应联络。适应联络本身并不难构造：规定 \(\nabla e_i=0\) 就保持了 \(J\)，但其挠率满足 \(T^\nabla(e_1,e_4)=-e_2\)。换用其他适应联络能够改变某些分量，却不能消去上面算出的 \((0,1)\) 投影。
\end{example}

最后回到\subsecref{g1:subsec:1305:04}的三维接触分布。设 \(D\subset\mathrm{T}M\) 为光滑秩二分布，令 \(V_0=\operatorname{span}\{\varepsilon_1,\varepsilon_2\}\subset\mathbb R^3\)。保留此标准平面的群为
\[
	G_D=\left\{\begin{pmatrix}B&b\\0&c\end{pmatrix}:
	B\in\operatorname{GL}(2,\mathbb R),\ b\in\mathbb R^2,\ c\neq0\right\}.
\]
局部取 \(D\) 的两个光滑基截面，再补一个横向向量，得到适应标架。因此
\[
	P_D=\{u\in\operatorname{Fr}(\mathrm{T}M):u(V_0)=D_{\pi(u)}\}
\]
是光滑 \(G_D\)-结构。其适应联络恰为保持分布的联络，即 \(Y\in\Gamma(D)\) 时，对任意 \(X\) 都有 \(\nabla_XY\in\Gamma(D)\)；这可由局部联络矩阵的左下块为零直接读出。

\ProseParagraphBreak

前文已经定义括号张量
\[
	\mathcal B_x(v,w)=[X,Y]_x\pmod{D_x},
	\qquad v,w\in D_x,\quad X_x=v,\ Y_x=w,
\]
其中 \(X,Y\) 为局部 \(D\) 截面。其值与延拓无关，因为 \([fX,Y]=f[X,Y]-Y(f)X\) 中的导数项仍属于 \(D\)。它是 \(\Lambda^2D^*\otimes(\mathrm{T}M/D)\) 的光滑截面，接触条件要求它在每点非零。对保持 \(D\) 的任意联络，挠率公式给出
\begin{equation}
	\label{g1:eq:distribution-intrinsic-torsion}
	T^\nabla(X,Y)\pmod D=-\mathcal B(X,Y),\qquad X,Y\in\Gamma(D),
\end{equation}
因为两项协变导数都在 \(D\) 中。从 \(\delta\) 也能看出同一事实：若 \(v,w\in V_0\)，而 \(A\) 取值于保持 \(V_0\) 的Lie 代数，则 \(\delta A(v,w)\in V_0\)，其在 \(V/V_0\) 中的分量为零。因此 \(-\mathcal B\) 是这个约化内蕴挠率的一个固定分量；接触分布的适应联络必定有挠。

\begin{example}[标准接触模型的挠率]
	\label{g1:exm:contact-intrinsic-torsion}
	沿用第十三章的
	\[
		\alpha=\mathrm{d}z-x\,\mathrm{d}y,\qquad
		A=\partial_x,\quad B=\partial_y+x\partial_z,
	\]
	并记 \(C=\partial_z\)。\((A,B,C)\) 是 \(D=\ker\alpha\) 的适应标架，对偶余标架为 \((\mathrm{d}x,\mathrm{d}y,\alpha)\)。\cref{g1:prop:contact-model-nonintegrable}已经证明 \([A,B]=C\)；故 \(\mathcal B(A,B)=[C]\neq0\)。令这个标架平行，便得到保持 \(D\) 的联络，其局部联络矩阵为零。因此第一结构方程给出挠率的分量列
	\[
		\begin{pmatrix}0\\0\\-\mathrm{d}x\wedge\mathrm{d}y\end{pmatrix},
		\qquad T^\nabla(A,B)=-C,
	\]
	与\eqref{g1:eq:distribution-intrinsic-torsion}完全相符。同时
	\[
		\mathrm{d}\alpha=-\mathrm{d}x\wedge\mathrm{d}y,
		\qquad
		\alpha\wedge\mathrm{d}\alpha
		=-\mathrm{d}x\wedge\mathrm{d}y\wedge\mathrm{d}z\neq0.
	\]
	这个已经写成标准形式的接触模型仍具有非零内蕴挠率。它的平面分布不能误认为坐标水平面 \(\ker\mathrm{d}z\)，后者的括号张量为零。
\end{example}

还要区分固定分布与固定定义形式。如果只要求 \(\ker\alpha=D\)，就允许将 \(\alpha\) 换成任意处处非零光滑倍数 \(f\alpha\)，上述群中的 \(c\) 可以取任意非零值。若把一个具体的 \(\alpha\) 也作为结构的一部分，并要求 \(u^*\alpha=\varepsilon^3\)，则适应标架间的变换必须满足 \(c=1\)。两者允许的结构群不同，因而允许的联络修正也不同。判断内蕴挠率之前，必须先明确几何结构究竟保留哪些数据。

\section{Cartan 联络与齐性模型}
\label{g1:sec:D02}

Euclid 平面上一副正向正交标架既有位置，也有方向。沿曲线移动这副标架时，需要同时记录位移与转动；只记录转动的普通主联络，没有把两类信息合在一起。Cartan 联络所做的事情，是把标架总空间在每一点的全部速度，与一个模型变换群的Lie 代数逐点对应。模型可以是Euclid 平面，也可以是球面或其他齐性空间；曲率衡量这种逐点对应能否在邻域中积分成与模型的实际等价。

\subsection{Cartan 联络的公理}
\label{g1:subsec:D02:01}

设 \(L\) 是有限维实Lie 群，\(H\subseteq L\) 是闭子群，Lie 代数分别为 \(\mathfrak l,\mathfrak h\)。由\cref{g1:thm:homogeneous-space}，商 \(L/H\) 是光滑流形，\(q:L\to L/H\) 是主右 \(H\) 丛，并在每个点附近有光滑局部截面。这里 \(L\) 描述模型的运动，\(H\) 描述固定模型基点 \(o=eH\) 的运动；二者的作用不同。以下流形均无边界、Hausdorff 且第二可数。

\ProseParagraphBreak

先看模型总空间 \(L\) 自身。它的左 Maurer--Cartan 形式为
\[
	(\theta_L)_a=(\mathrm{d}L_{a^{-1}})_a:
	\mathrm{T}_aL\longrightarrow\mathfrak l.
\]
它用左平移把每个速度搬回单位元，是逐点线性同构。若只把 \(L\) 看成主 \(H\) 丛，垂直速度为
\[
	X_a^\#=\left.\frac{\mathrm{d}}{\mathrm{d}t}\right|_{t=0}
	a\exp(tX),\qquad X\in\mathfrak h,
\]
所以 \(\theta_L(X^\#)=X\)。\cref{g1:prop:maurer-cartan-product}给出的乘法公式又说明 \(R_h^*\theta_L=\operatorname{Ad}_{h^{-1}}\theta_L\)。把这三项性质保留下来，就能在一个未必齐性的空间上规定相同类型的无穷小模型。

\begin{definition}[Cartan 联络]\label{g1:def:cartan-connection}
	设 \(L\) 为有限维实Lie 群，\(H\subseteq L\) 为闭子群，\(\pi:P\to M\) 为无边界光滑流形上的光滑主右 \(H\) 丛。若 \(\mathfrak l\) 值光滑一形式 \(\omega\in\Omega^1(P;\mathfrak l)\) 满足下列条件，就称它为以 \(L/H\) 为模型的\term[cartanlianluo]{Cartan 联络}[Cartan connection]：
	\begin{enumerate}
		\item 对每个 \(p\in P\)，\(\omega_p:\mathrm{T}_pP\to\mathfrak l\) 是线性同构；
		\item 对每个 \(h\in H\)，有 \(R_h^*\omega=\operatorname{Ad}_{h^{-1}}\omega\)；
		\item 对每个 \(X\in\mathfrak h\)，有 \(\omega(X^\#)=X\)，其中 \(X^\#_p=\left.\frac{\mathrm{d}}{\mathrm{d}t}\right|_0p\exp(tX)\)。
	\end{enumerate}
	二元组 \((P,\omega)\) 称为一个模型为 \(L/H\) 的\term[cartanjihe]{Cartan 几何}[Cartan geometry]。两个 Cartan 几何之间的同构是覆盖底空间微分同胚的主 \(H\) 丛同构；沿它拉回目标的 Cartan 联络，应当得到源的 Cartan 联络。
\end{definition}

第一条公理称为\term[jueduipingxing]{绝对平行}[absolute parallelism]条件：固定 \(\mathfrak l\) 的一组基，便得到 \(P\) 上一组处处线性无关的向量场。这个名称不表示这些向量场交换，也不表示某个已经给定的协变导数使它们平行。第二条说明改变标架后的速度记录如何变化，第三条则把纯粹的标架变化与 \(\mathfrak h\) 中的元素对应。三条公理同时满足时，沿底空间运动的信息也已经包含在 \(\omega\) 中。

\begin{proposition}[Cartan 联络的焊接同构]\label{g1:prop:cartan-soldering}
	设 \((P\to M,\omega)\) 为模型 \(L/H\) 的光滑 Cartan 几何，\(L\) 为有限维实Lie 群，\(H\) 为闭子群。则
	\[
		\dim M=\dim L-\dim H=\dim(L/H),\qquad
		\omega_p^{-1}(\mathfrak h)=V_pP.
	\]
	商表示 \(h\cdot(X+\mathfrak h)=\operatorname{Ad}_hX+\mathfrak h\) 给出自然的向量丛同构
	\begin{equation}\label{g1:eq:cartan-soldering}
		\mathrm{T}M\cong P\times_H(\mathfrak l/\mathfrak h).
	\end{equation}
\end{proposition}
\begin{proof}
	绝对平行条件给出 \(\dim P=\dim\mathfrak l\)，主丛局部乘积给出 \(\dim P=\dim M+\dim H\)，从而得到维数等式。由\cref{g1:prop:principal-fundamental-fields}，\(X\mapsto X_p^\#\) 是 \(\mathfrak h\to V_pP\) 的同构；再现条件说明 \(\omega_p\) 在这两个子空间之间正是它的逆。由于 \(\omega_p\) 在整个切空间上也单射，\(\mathfrak h\) 的原像恰为 \(V_pP\)。

	给定 \(v\in\mathrm{T}_xM\) 和 \(p\in P_x\)，选取 \(\widetilde v\in\mathrm{T}_pP\) 满足 \((\mathrm{d}\pi)_p\widetilde v=v\)，定义
	\[
		\overline\omega_p(v)=\omega_p(\widetilde v)+\mathfrak h.
	\]
	两个提升之差垂直，故其 \(\omega\) 值之差在 \(\mathfrak h\) 中，这一定义与提升无关。前段的原像等式说明 \(\overline\omega_p\) 单射；维数相同，故它是同构。

	因为 \(\operatorname{Ad}_h\mathfrak h=\mathfrak h\)，商表示有定义。等变性给出
	\[
		\overline\omega_{ph}(v)
		=\operatorname{Ad}_{h^{-1}}\overline\omega_p(v).
	\]
	关联丛约定为 \([ph,\xi]=[p,h\cdot\xi]\)，见\cref{g1:thm:associated-vector-bundle}，所以
	\[
		v\longmapsto[p,\overline\omega_p(v)]
	\]
	与 \(p\) 的选择无关。它在局部截面中由光滑可逆矩阵给出，因此是所求光滑向量丛同构。
\end{proof}

这里的商值一形式把抽象模型方向 \(\mathfrak l/\mathfrak h\) 与真实切向量联系起来，因而称为焊接形式。它没有要求先选一个补空间 \(\mathfrak l=\mathfrak h\oplus\mathfrak m\)；商空间本身已经足够。特别地，Cartan 联络不是任意主 \(H\) 丛上都能附加的数据，底空间维数和切丛的关联表示必须相容。

\ProseParagraphBreak

为了计算，取局部截面 \(s:U\to P\)，记 \(B=s^*\omega\)，并用 \(\Psi(x,h)=s(x)h\) 平凡化。与\cref{g1:thm:principal-connection-local-data}中的推导一样，把 \(\mathrm{d}\Psi(v,\eta)\) 拆成 \((\mathrm{d}R_h)\mathrm{d}s(v)\) 和垂直部分；分别用等变性与再现条件，便得到
\begin{equation}\label{g1:eq:cartan-local-form}
	\Psi^*\omega
	=\operatorname{Ad}_{h^{-1}}B+\theta_H,
	\qquad
	B\bmod\mathfrak h:
	\mathrm{T}_xU\longrightarrow\mathfrak l/\mathfrak h
	\text{ 为同构}.
\end{equation}
这里 \(\theta_H\) 是 \(H\) 的左 Maurer--Cartan 形式，通过包含映射 \(\mathfrak h\subseteq\mathfrak l\)，将其视为 \(\mathfrak l\) 值形式；第一项作用于 \(U\) 方向，第二项作用于 \(H\) 方向。若换成 \(s'=s\,k\)，其中 \(k:U\to H\) 光滑，则
\[
	B'=\operatorname{Ad}_{k^{-1}}B+k^*\theta_H.
\]
这个变换律与普通主联络相似，但 \(B\) 还含有不能丢掉的 \(\mathfrak l/\mathfrak h\) 分量。

\subsection{曲率、平坦性及路径展开}
\label{g1:subsec:D02:02}

模型上的 \(\theta_L\) 不只是逐点可逆，它还满足一个微分相容条件。对左不变场 \(X^L,Y^L\)，有 \(\theta_L(X^L)=X\)、\(\theta_L(Y^L)=Y\)，二者为常值，因此
\[
	\mathrm{d}\theta_L(X^L,Y^L)
	=-\theta_L([X^L,Y^L])=-[X,Y].
\]
左不变场逐点张成切空间，故
\begin{equation}\label{g1:eq:cartan-model-mc}
	\mathrm{d}\theta_L+\tfrac12[\theta_L\wedge\theta_L]=0.
\end{equation}
这里沿用\eqref{g1:eq:lie-bracket-wedge}的约定，即
\[
	[\omega\wedge\omega](U,V)=2[\omega(U),\omega(V)].
\]
因子 \(1/2\) 恰好抵消交替求和产生的重复。

\begin{definition}\label{g1:def:cartan-curvature}
	设 \((P\to M,\omega)\) 为模型 \(L/H\) 的光滑 Cartan 几何。定义它的\term[cartanqulv]{Cartan 曲率}[Cartan curvature]为
	\begin{equation}\label{g1:eq:cartan-curvature}
		\Omega=\mathrm{d}\omega+\tfrac12[\omega\wedge\omega]
		\in\Omega^2(P;\mathfrak l).
	\end{equation}
	若 \(\Omega=0\)，称这个 Cartan 几何平坦。
\end{definition}

\begin{proposition}[曲率的水平性、等变性与 Bianchi 恒等式]\label{g1:prop:cartan-curvature-identities}
	设 \((P\to M,\omega)\) 为模型 \(L/H\) 的光滑 Cartan 几何，其中 \(L\) 为有限维实Lie 群、\(H\subseteq L\) 闭。其曲率满足
	\begin{equation}\label{g1:eq:cartan-curvature-properties}
		\iota_{X^\#}\Omega=0\quad(X\in\mathfrak h),\qquad
		R_h^*\Omega=\operatorname{Ad}_{h^{-1}}\Omega\quad(h\in H),
	\end{equation}
	以及
	\begin{equation}\label{g1:eq:cartan-bianchi}
		\mathrm{d}\Omega+[\omega\wedge\Omega]=0.
	\end{equation}
\end{proposition}
\begin{proof}
	将 \(R_{\exp(tX)}^*\omega=\operatorname{Ad}_{\exp(-tX)}\omega\) 在零点求导，得
	\[
		\mathcal L_{X^\#}\omega=-[X,\omega].
	\]
	由于 \(\iota_{X^\#}\omega=X\) 是常函数，Cartan 的 Lie 导数公式给出
	\[
		\iota_{X^\#}\mathrm{d}\omega=-[X,\omega].
	\]
	另一方面，插入括号楔积给出
	\(\frac12\iota_{X^\#}[\omega\wedge\omega]=[X,\omega]\)，两项抵消，证明水平性。固定 \(h\in H\) 时，\(\operatorname{Ad}_{h^{-1}}\) 是常系数的Lie 代数自同构；它与外微分相交换并保持括号。因此直接拉回\eqref{g1:eq:cartan-curvature}即得等变性。

	Bianchi 恒等式只使用Lie 代数括号的代数性质，不要求 \(\omega\) 是 \(\mathfrak h\) 值。具体地，\cref{g1:thm:principal-bianchi}中已由分量计算得到
	\[
		\mathrm{d}[\omega\wedge\omega]
		=-2[\omega\wedge\mathrm{d}\omega],
		\qquad
		[\omega\wedge[\omega\wedge\omega]]=0.
	\]
	第二式是对三个输入交替求和后的 Jacobi 恒等式。于是
	\[
		\mathrm{d}\Omega=-[\omega\wedge\mathrm{d}\omega],
		\qquad
		[\omega\wedge\Omega]=[\omega\wedge\mathrm{d}\omega],
	\]
	相加便得结论。
\end{proof}

曲率的“水平性”表示有一个垂直输入时曲率为零，因此它通过底空间的切向量取值。结合焊接同构，可以把 \(\Omega_p\) 写成
\[
	\kappa(p)\in
	\operatorname{Hom}\bigl(\Lambda^2(\mathfrak l/\mathfrak h),\mathfrak l\bigr),
	\qquad
	\kappa(p)(\overline X,\overline Y)
	=\Omega_p(\omega_p^{-1}X,\omega_p^{-1}Y).
\]
改变 \(X,Y\) 的代表只会给输入增加垂直向量，所以这个式子有定义。它说明 Cartan 曲率记录的是两个底空间方向之间的相容障碍。

\begin{example}[定向曲面的Euclid  Cartan 几何]\label{g1:exm:cartan-euclidean-surface}
	设 \(M\) 为定向的二维光滑Riemann 流形，\(P\to M\) 为其正向正交标架丛。把标架 \(p=(e_1,e_2)\) 看作等距同构 \(p:\mathbb R^2\to\mathrm{T}_{\pi(p)}M\)，定义焊接形式
	\[
		\theta_p(V)=p^{-1}(\mathrm{d}\pi)_pV
		=\begin{pmatrix}\theta^1_p(V)\\\theta^2_p(V)\end{pmatrix}.
	\]
	它在垂直向量上为零，且 \(R_h^*\theta=h^{-1}\theta\)。令 \(A=J_0\alpha\) 为该度量的无挠度量联络在 \(P\) 上的主联络形式，其中
	\[
		J_0=\begin{pmatrix}0&-1\\1&0\end{pmatrix},
		\qquad
		\alpha\in\Omega^1(P).
	\]
	局部截面拉回的 \(\alpha\) 正是\cref{g1:def:surface-connection-form}中满足 \(\nabla e_1=\alpha e_2\)、\(\nabla e_2=-\alpha e_1\) 的形式；其换架公式把这些局部数据粘成上述主联络。

	取 \(L=\operatorname{SE}(2)=\mathbb R^2\rtimes\operatorname{SO}(2)\)，\(H=\operatorname{SO}(2)\)，于是 \(L/H\) 是定向Euclid 平面。\(\mathfrak l=\mathfrak{so}(2)\ltimes\mathbb R^2\) 的括号为
	\[
		[(B,u),(C,v)]=([B,C],Bv-Cu).
	\]
	在 \(P\) 上令
	\[
		\omega=(J_0\alpha,\theta).
	\]
	若 \(\omega_p(V)=0\)，则 \(\theta_p(V)=0\) 使 \(V\) 垂直，再现条件又使 \(A_p(V)=0\) 只能发生于 \(V=0\)。两边维数均为 \(3\)，故 \(\omega_p\) 是同构。主联络与焊接形式的等变律，恰好分别是 \(\operatorname{Ad}_{h^{-1}}\) 在旋转分量与平移分量上的作用；垂直再现也成立。因此 \(\omega\) 是 Cartan 联络。

	按半直积括号展开曲率，得到
	\begin{equation}\label{g1:eq:cartan-surface-curvature}
		\Omega=
		\left(
		J_0\,\mathrm{d}\alpha,\ 
		\mathrm{d}\theta+J_0\alpha\wedge\theta
		\right).
	\end{equation}
	第二分量是切丛联络的挠率形式。此处联络无挠，因而它为零；分量式是
	\[
		\mathrm{d}\theta^1=\alpha\wedge\theta^2,\qquad
		\mathrm{d}\theta^2=-\alpha\wedge\theta^1.
	\]
	由\cref{g1:prop:connection-curvature-form}，第一分量满足
	\[
		\mathrm{d}\alpha=-(K\circ\pi)\theta^1\wedge\theta^2.
	\]
	该式在局部截面上就是 \(\mathrm{d}\alpha=-K\,\mathrm{d}A\)；两边都水平且等变，所以它在整个标架丛上成立。因此这个以Euclid 平面为模型的 Cartan 几何平坦，当且仅当 \(K=0\)。单位球面虽然无挠，仍有非零的Euclid  Cartan 曲率。
\end{example}

为什么要把旋转曲率与挠率放进同一个 \(\Omega\)？局部Euclid 坐标不仅要求标架在各方向上的转动相容，也要求记录的位移真正来自坐标。前者对应旋转分量，后者对应平移分量。对一般齐性模型，下面的定理把这两种相容要求统一为 Maurer--Cartan 方程。

\begin{theorem}[平坦 Cartan 几何的局部模型]\label{g1:thm:flat-cartan-local-model}
	设 \(L\) 为有限维实Lie 群，\(H\subseteq L\) 为闭子群，\((P\to M,\omega)\) 为无边界光滑流形上的 Cartan 几何。记 \(q:L\to L/H\) 为商投影，\((\theta_L)_a=(\mathrm dL_{a^{-1}})_a\) 为左 Maurer--Cartan 形式。其曲率为零，当且仅当每个 \(x_0\in M\) 都有开邻域 \(U\)，以及 \(L/H\) 中的开集 \(V\)，使 \((P|_U,\omega)\) 与 \((q^{-1}(V),\theta_L)\) 作为主 \(H\) 丛及 Cartan 联络同构。
\end{theorem}
\begin{proof}
	先设 \(\Omega=0\)。我们先在总空间的一点附近积分 \(\omega\)，再借助主丛局部截面将所得映射扩展到整条 \(H\) 纤维。第二步是从总空间的局部坐标得到底空间局部等价所必需的。

	对 \(X\in\mathfrak l\)，令 \(E_X\) 为由 \(\omega(E_X)=X\) 定义的光滑向量场。由曲率为零，且等式两边均为常值，可得
	\[
		0=\Omega(E_X,E_Y)
		=-\omega([E_X,E_Y])+[X,Y],
		\qquad [E_X,E_Y]=E_{[X,Y]}.
	\]
	在 \(P\times L\) 上考虑分布
	\[
		\mathcal D_{(p,a)}
		=\{(v,w):\omega_p(v)=(\theta_L)_a(w)\}.
	\]
	两形式都逐点可逆，故此分布秩恒为 \(\dim L\)。若 \(X_i\) 是 \(\mathfrak l\) 的一组基，则 \((E_{X_i},X_i^L)\) 给出它的光滑标架，且
	\[
		[(E_X,X^L),(E_Y,Y^L)]
		=(E_{[X,Y]},[X,Y]^L).
	\]
	标架的括号仍在分布中；再用带函数系数的括号公式，得 \(\mathcal D\) 对合。

	固定 \(p_0\in P_{x_0}\)。\cref{g1:thm:frobenius}适用于无边界流形 \(P\times L\) 上的这个恒秩光滑分布，给出通过 \((p_0,e)\) 的局部嵌入积分片。第一投影在 \(\mathcal D\) 上的微分为同构，因为每个 \(v\) 唯一决定 \(w\)。由\cref{g1:thm:inverse-function}，缩小积分片后，它是某个光滑映射 \(f:W\to L\) 的图像，其中 \(W\) 为 \(p_0\) 的开邻域、\(f(p_0)=e\)。图像的切空间在 \(\mathcal D\) 中，正好说明
	\[
		f^*\theta_L=\omega|_W.
	\]
	两个一形式逐点可逆，故 \(\mathrm{d}f\) 逐点可逆。

	取经过 \(p_0\) 的局部截面 \(s:U\to P\)，缩小 \(U\) 使 \(s(U)\subset W\)，并记 \(a=f\circ s\)。定义
	\[
		\Phi:P|_U\longrightarrow L,\qquad
		\Phi(s(x)h)=a(x)h,\quad h\in H.
	\]
	主丛平凡化给出 \(s(x)h\) 的唯一表示，故 \(\Phi\) 光滑、有定义，并对所有 \(h\in H\) 等变；这里不要求 \(H\) 连通。沿截面有 \(a^*\theta_L=s^*\omega\)。由 Maurer--Cartan 乘法公式及\eqref{g1:eq:cartan-local-form}，在 \(U\times H\) 上
	\[
		\Phi^*\theta_L
		=\operatorname{Ad}_{h^{-1}}a^*\theta_L+\theta_H
		=\operatorname{Ad}_{h^{-1}}s^*\omega+\theta_H
		=\omega.
	\]
	此式是在主丛平凡化中书写的，因而给出整个 \(P|_U\) 上的等式。

	定义底映射 \(\phi(x)=a(x)H\)。它的微分可由 \(s^*\omega\bmod\mathfrak h\) 计算：左平移 \(a(x)^{-1}\) 将 \(\mathrm{T}_{\phi(x)}(L/H)\) 识别为 \(\mathfrak l/\mathfrak h\)，在此识别下，\(\mathrm{d}\phi_x\) 就是 \(\overline\omega_{s(x)}\)。\cref{g1:prop:cartan-soldering}说明它为同构。再缩小 \(U\)，逆函数定理使 \(\phi:U\to V\) 成为到某个开集 \(V\subseteq L/H\) 的微分同胚。

	对每个 \(x\in U\)，\(\Phi\) 将整条纤维 \(s(x)H\) 双射到 \(a(x)H=q^{-1}(\phi(x))\)。因此它是 \(P|_U\to q^{-1}(V)\) 的双射。目标截面 \(y\mapsto a(\phi^{-1}(y))\) 给出局部平凡化，在两边平凡化中 \(\Phi\) 写成 \((x,h)\mapsto(\phi(x),h)\)，故其逆光滑。这就得到保持 Cartan 联络的主丛同构。

	反之，若每一点附近存在这样的同构，拉回模型上的\eqref{g1:eq:cartan-model-mc}，便得 \(\Omega|_{P|_U}=0\)。这些 \(U\) 覆盖 \(M\)，故 \(\Omega=0\)。
\end{proof}

证明中的积分片把两个空间里由同一个 \(X\in\mathfrak l\) 指定的运动配对；曲率为零恰好保证这些配对在取括号后仍相容。局部截面随后把积分片产生的映射延伸到整条 \(H\) 纤维，才得到底空间开片上的几何等价。

\ProseParagraphBreak

曲率不为零时，仍然可以沿一条指定路径，把实际空间中的运动记录到模型中。给定有限分片光滑路径 \(\gamma:[0,1]\to M\) 和初始标架 \(p_0\in P_{\gamma(0)}\)，先选取同样正则的提升 \(p(t)\)，满足 \(p(0)=p_0\)。这种提升可在有限个局部平凡化中逐段构造：每段保持选定的群坐标不变，在接点调节后一段的群坐标使两段接合。

\begin{proposition}[路径的展开]\label{g1:prop:cartan-path-development}
	设 \((P\to M,\omega)\) 为模型 \(L/H\) 的光滑 Cartan 几何，\(L\) 为有限维实Lie 群，\(H\subseteq L\) 闭。给定有限分片光滑路径 \(\gamma:[0,1]\to M\) 及 \(p_0\in P_{\gamma(0)}\)，选择有限分片光滑提升 \(p:[0,1]\to P\)，使 \(\pi\circ p=\gamma\)、\(p(0)=p_0\)。记 \((\theta_L)_a=(\mathrm dL_{a^{-1}})_a\) 为 \(L\) 上的左 Maurer--Cartan 形式。则未知路径 \(g:[0,1]\to L\) 的初值问题
	\begin{equation}\label{g1:eq:cartan-development-ode}
		\theta_L(g'(t))=\omega(p'(t)),\qquad g(0)=e
	\end{equation}
	在整个区间上有唯一有限分片光滑解。路径
	\[
		\operatorname{dev}_{p_0}(\gamma)(t)=g(t)H
	\]
	与保持 \(p(0)=p_0\) 的提升选择无关，称为 \(\gamma\) 相对于 \(p_0\) 的\term[lujingzhankai]{路径展开}[development]。对任意 \(h_0\in H\)，若把初始标架改成 \(p_0h_0\)，则展开路径变为 \(h_0^{-1}g(t)H\)。
\end{proposition}
\begin{proof}
	记 \(b(t)=\omega(p'(t))\)。方程等价于 \(g'=(\mathrm{d}L_g)_e b(t)\)。令 \(u=g^{-1}\)，对 \(g(t)u(t)=e\) 求导可得
	\[
		u'=-(\mathrm{d}R_u)_e b(t).
	\]
	\cref{g1:lem:right-invariant-time-ode}说明，光滑Lie 代数值系数在紧时间区间上定义的这类右平移方程有全区间唯一解；其证明用单位元附近的统一短时存在性及右平移延拓，不要求 \(L\) 紧。对每个光滑子区间应用该结论并接合初值，即得此处的解。

	另一个提升可唯一写成 \(\widehat p(t)=p(t)h(t)\)，其中 \(h(t)\in H\) 有限分片光滑。对变化的 \(h(t)\)，等变性和再现条件给出
	\[
		\omega(\widehat p')
		=\operatorname{Ad}_{h^{-1}}\omega(p')+\theta_H(h').
	\]
	另一方面，\cref{g1:prop:maurer-cartan-product}给出
	\[
		\theta_L((gh)')
		=\operatorname{Ad}_{h^{-1}}\theta_L(g')+\theta_H(h').
	\]
	若两提升初值均为 \(p_0\)，则 \(h(0)=e\)，所以 \(gh\) 满足新方程及单位元初值。唯一性使新解正是 \(gh\)，投影后 \((gh)H=gH\)，从而展开路径不变。

	若初始标架改为 \(p_0h_0\)，可选择提升 \(ph_0\)。新系数为 \(\operatorname{Ad}_{h_0^{-1}}b\)，新解为 \(h_0^{-1}gh_0\)，因为它在零时刻为单位元且满足相应方程。其陪集为 \(h_0^{-1}gH\)，得到最后的变换式。
\end{proof}

在\cref{g1:thm:flat-cartan-local-model}的平坦邻域中，若局部同构归一化为 \(\Phi(p_0)=e\)，则 \(\Phi(p(t))\) 满足展开方程，故展开就是底空间局部映射 \(\phi(\gamma(t))\)。一般情形下，展开可以依赖整条路径；不能仅由终点 \(x\) 写出一个单值的 \(\operatorname{dev}(x)\)。即使曲率为零，把不同邻域的映射接成全局单值映射也还涉及绕环路的延拓，\secref{g1:sec:D03}将对此作具体计算。

\ProseParagraphBreak

Cartan 几何的完备性提出反向的延拓问题：给定模型中的运动，实际空间中是否始终存在相应运动？这与沿实际空间中已给路径求解展开方程使用的数据不同。例如开区间 \((0,1)\) 上的 \(\omega=\mathrm{d}x\) 以加法群 \(\mathbb R\) 为模型，曲率为零；但从 \(x=1/2\) 按单位速度前进，只能在到达区间边界之前留在原空间中。局部平坦、全局单值和运动能否无限延拓，是三个不同问题。

\subsection{Cartan 联络与主联络}
\label{g1:subsec:D02:03}

第18章通过一个与右作用相容的水平补空间定义主联络，这也是主丛上的 Ehresmann 联络。\footnote{Charles Ehresmann（1905--1979），法国数学家，研究纤维丛、联络与微分拓扑。}若用一形式表示，它是 \(\mathfrak h\) 值的 \(\eta\)，满足等变性和垂直再现，并由 \(\ker\eta\) 给出水平分布。Cartan 联络则是更大Lie 代数 \(\mathfrak l\) 值的逐点同构，因此
\[
	\ker\omega_p=0,\qquad
	\omega_p^{-1}(\mathfrak h)=V_pP.
\]
把 \(\ker\omega\) 当成水平分布，会使所有非零底空间速度都无从提升。

\ProseParagraphBreak

两类形式在垂直方向上的行为相同，但在底空间方向上承担的任务不同。\cref{g1:tab:cartan-ehresmann}把这一区别列在同一个主 \(H\) 丛上；其中 \(\mathcal H\) 表示水平分布，避免与结构群 \(H\) 混用。
\begin{table}[htbp]
	\centering
	\caption{主联络与 Cartan 联络所记录的数据}
	\label{g1:tab:cartan-ehresmann}
	\begin{tabularx}{\linewidth}{@{}lXX@{}}
		\toprule
		比较项 & 主丛 Ehresmann 联络 \(\eta\) & 模型 \(L/H\) 的 Cartan 联络 \(\omega\)\\
		\midrule
		取值空间 & \(\mathfrak h\) & \(\mathfrak l\)\\
		逐点作用 & 记录垂直分量，\(\eta|_{VP}\) 可逆 & 记录全部速度，\(\omega_p\) 可逆\\
		核 & \(\mathcal H_p\)，维数为 \(\dim M\) & 零空间\\
		底空间方向 & 投影 \(\mathcal H_p\to\mathrm{T}_xM\) 为同构 & 选定 \(p\in P_x\) 后给出 \(\mathrm{T}_xM\cong\mathfrak l/\mathfrak h\)\\
		维数限制 & 与某个外部模型群无关 & \(\dim M=\dim(L/H)\)\\
		平坦性的局部含义 & 每个总空间点均有局部水平截面通过 & 与 \(L/H\) 的开片局部等价\\
		\bottomrule
	\end{tabularx}
\end{table}

若模型Lie 代数存在与 \(H\) 作用相容的补空间，两类联络仍有直接联系。这里必须要求补空间在整个 \(\operatorname{Ad}H\) 下不变；任取一个线性补空间虽能分解向量，却未必能分解出满足主丛等变律的形式。

\begin{proposition}[不变补空间下的分解]\label{g1:prop:reductive-cartan-splitting}
	设 \(L\) 为有限维实Lie 群，\(H\subseteq L\) 为闭子群，且有 \(\operatorname{Ad}H\) 不变的直和
	\[
		\mathfrak l=\mathfrak h\oplus\mathfrak m.
	\]
	对任意以 \(L/H\) 为模型的光滑 Cartan 几何 \((P\to M,\omega)\)，记
	\[
		\eta=\operatorname{pr}_{\mathfrak h}\omega,\qquad
		\vartheta=\operatorname{pr}_{\mathfrak m}\omega.
	\]
	则 \(\eta\) 是主 \(H\) 联络，\(\vartheta\) 为水平、等变的 \(\mathfrak m\) 值形式，且诱导 \(\mathrm{T}M\cong P\times_H\mathfrak m\)。此时主联络的水平空间为 \(\omega^{-1}(\mathfrak m)\)。
\end{proposition}
\begin{proof}
	两直和分量都在 \(\operatorname{Ad}H\) 下不变，故两投影与 \(\operatorname{Ad}H\) 交换。\(\omega\) 的等变性因此分别下降为 \(\eta,\vartheta\) 的等变性。对 \(X\in\mathfrak h\)，再现条件给出 \(\eta(X^\#)=X\)、\(\vartheta(X^\#)=0\)。由\cref{g1:thm:principal-connection-equivalence}，前两项性质使 \(\eta\) 成为主联络，而后一等式给出 \(\vartheta\) 的水平性。投影 \(\mathfrak m\to\mathfrak l/\mathfrak h\) 是 \(H\) 等变同构，故焊接同构直接成为 \(\mathrm{T}M\cong P\times_H\mathfrak m\)。最后，\(\eta_p(v)=0\) 等价于 \(\omega_p(v)\in\mathfrak m\)，所以得到所述水平空间。
\end{proof}

以上讨论的是具有不变补空间的模型（reductive model）。即使不变补空间存在，也可能不唯一；由它投影得到的普通主联络随所选分解而定。如果不存在不变补空间，就不能用一个任意线性投影，从 \(\omega\) 自然得到普通主联络。主 \(H\) 丛本身仍可另选主联络，这是另一项选择。

\ProseParagraphBreak

在这个分解下，Cartan 曲率的两分量写成
\[
	\begin{aligned}
		\Omega_{\mathfrak h}
		&=\mathrm{d}\eta+\tfrac12[\eta\wedge\eta]
		+\tfrac12[\vartheta\wedge\vartheta]_{\mathfrak h},\\
		\Omega_{\mathfrak m}
		&=\mathrm{d}\vartheta+[\eta\wedge\vartheta]
		+\tfrac12[\vartheta\wedge\vartheta]_{\mathfrak m}.
	\end{aligned}
\]
这里 \([\mathfrak h,\mathfrak m]\subseteq\mathfrak m\) 来自不变性，因此混合项只出现在第二行。Euclid 模型的平移部分满足 \([\mathfrak m,\mathfrak m]=0\)，于是两行分别是普通联络的曲率与挠率；\cref{g1:exm:cartan-euclidean-surface}正是这种情形。一般模型中还要保留 \([\vartheta\wedge\vartheta]\) 的贡献，所以“Cartan 平坦”必须连同所选模型一起理解。

\section{标架归一化与几何等价}
\label{g1:sec:D03}

同一张平面可以用直角标架或旋转标架描述，联络系数随之改变；周长不同的平圆柱却可以具有相同的零曲率，而不具有相同的整体几何。等价问题必须同时回答两件事：哪些差异只是表示方式的选择，哪些数据在改变表示后仍然存在。本节先在结构方程中辨认这些自由，再用曲率不恒定的度量曲面给出一次具体的标架归一化，最后用显式展开映射检验局部结论能否成为整体结论。

\subsection{结构方程中的规范自由}
\label{g1:subsec:D03:01}

设 \(G\subseteq\operatorname{GL}(n,\mathbb R)\) 是矩阵Lie 群，在一个带 \(G\)-结构的 \(n\) 维流形的开集 \(U\) 上选适应标架 \(e=(e_1,\ldots,e_n)\)。对偶余标架写成列 \(\theta=(\theta^1,\ldots,\theta^n)^{\mathsf T}\)，即 \(V=e\,\theta(V)\)。选定一个相容联络后，沿用 \(\nabla e=eA\)，其中 \(A\) 是取值于 \(\mathfrak g\) 的局部联络矩阵。

\begin{proposition}[换标架与结构方程]\label{g1:prop:equivalence-frame-change}
	设 \(G\subseteq\operatorname{GL}(n,\mathbb R)\) 为矩阵Lie 群，\(e\) 是某个 \(n\) 维光滑流形的开集 \(U\) 上一个 \(G\)-结构的光滑适应标架，\(\nabla\) 为相容联络，\(\theta,A\) 分别为此标架中的余标架与联络矩阵。将标架改为 \(e'=eg\)，其中 \(g:U\to G\) 光滑，联络本身保持不变。则
	\begin{equation}\label{g1:eq:equivalence-frame-change}
		\theta'=g^{-1}\theta,\qquad
		A'=g^{-1}Ag+g^{-1}\mathrm{d}g.
	\end{equation}
	记挠率形式及曲率矩阵为 \(T=\mathrm{d}\theta+A\wedge\theta\)、\(F_A=\mathrm{d}A+A\wedge A\)，则
	\[
		T'=g^{-1}T,\qquad F_{A'}=g^{-1}F_Ag.
	\]
	这里矩阵值形式的楔积同时使用矩阵乘法，向量值形式的楔积按矩阵作用在列向量上解释。
\end{proposition}
\begin{proof}
	恒等式 \(e\theta(V)=eg\theta'(V)\) 对每个切向量成立，所以 \(\theta'=g^{-1}\theta\)。再对 \(e'=eg\) 使用联络的乘法法则，得到
	\[
		\nabla e'=e(Ag+\mathrm{d}g)=e'\bigl(g^{-1}Ag+g^{-1}\mathrm{d}g\bigr),
	\]
	即第二式。这也恢复了\cref{g1:thm:principal-connection-local-data}在标架丛上的局部公式：固定的主联络在不同截面上的表达不同。

	由 \(\mathrm{d}(g^{-1})=-g^{-1}(\mathrm{d}g)g^{-1}\)，
	\[
		\begin{aligned}
			\mathrm{d}\theta'&=-g^{-1}\mathrm{d}g\wedge g^{-1}\theta+g^{-1}\mathrm{d}\theta,\\
			A'\wedge\theta'&=g^{-1}A\wedge\theta+g^{-1}\mathrm{d}g\wedge g^{-1}\theta.
		\end{aligned}
	\]
	相加时含 \(\mathrm{d}g\) 的两项抵消，得到挠率公式。对联络的算子 \(\mathrm{D}=\mathrm{d}+A\)，新分量中的算子为 \(\mathrm{D}'=g^{-1}\mathrm{D}g\)；连续作用两次得到 \((\mathrm{D}')^2=g^{-1}\mathrm{D}^2g\)。由于 \(\mathrm{D}^2\) 在列向量值函数上等于左乘 \(F_A\)，曲率公式随之成立。
\end{proof}

这里的换标架没有移动底空间的点，也没有改变联络。它与\cref{g1:prop:gauge-connection-formula}中的主动规范变换具有同样的局部公式，但主动规范变换将联络拉回到另一个联络，所选截面可以保持不动。另一方面，两张 \(G\)-结构之间的几何等价还要有底空间的局部微分同胚，并要求其微分把适应标架送到适应标架。这是对底点和切向量的共同要求，不能由任取一族矩阵 \(g(x)\) 代替。

\ProseParagraphBreak

例如，设 \(\Phi:U\to\widehat U\) 是待检验的局部微分同胚，两边的标架分别为 \(e,\widehat e\)。若写
\[
	\Phi_*e=\widehat e\,k,\qquad k:U\to G,
\]
则余标架必须满足 \(\Phi^*\widehat\theta=k\theta\)。若还要求保持所选联络，则必须同时满足
\[
	A=k^{-1}(\Phi^*\widehat A)k+k^{-1}\mathrm{d}k.
\]
第二式由对 \(\Phi_*e=\widehat e k\) 求协变导数得到，第一式则直接来自切向量的分量。两个条件中的同一个 \(k\) 必须与 \(\mathrm{d}\Phi\) 相容。只把联络矩阵变成某种相似形式，尚未得到这样的 \(\Phi\)。

\ProseParagraphBreak

另一个不同的选择是固定标架，改变相容联络。若 \(A\) 改为 \(A+B\)，其中 \(B\in\Omega^1(U;\mathfrak g)\)，则
\[
	T_{A+B}=T_A+B\wedge\theta.
\]
用 \(\theta\) 将切空间识别为 \(V=\mathbb R^n\)，在一点把 \(B\) 看作 \(V^*\otimes\mathfrak g\) 的元素，挠率的变化就是
\[
	(\delta B)(v,w)=B(v)w-B(w)v.
\]
这正是\subsecref{g1:subsec:D01:02}中讨论的线性映射。给定挠率中落在 \(\operatorname{im}\delta\) 的一项，才可能选 \(B\) 将它消去；商空间中的内蕴挠率不受这一操作影响。如果有多个 \(B\) 能消去同一项，它们之差落在 \(\ker\delta\)。因此，选联络消去挠率与选标架减少结构群是两种操作，各自留下的自由也不同。

\begin{example}[二维正交结构中的消挠率]\label{g1:exm:surface-torsion-absorption}
	在定向Riemann 度量曲面上取正向正交余标架 \(\theta^1,\theta^2\)，记 \(\Theta=\theta^1\wedge\theta^2\)。由于二维二形式只有一个独立分量，存在唯一的一对光滑函数 \(c_1,c_2\) 使
	\[
		\mathrm{d}\theta^1=c_1\Theta,\qquad
		\mathrm{d}\theta^2=c_2\Theta.
	\]
	所有度量相容的局部联络矩阵都形如
	\[
		A=J_0\alpha,\qquad
		J_0=\begin{pmatrix}0&-1\\1&0\end{pmatrix},\qquad
		\alpha=a\theta^1+b\theta^2.
	\]
	把它代入挠率公式，逐项得到
	\[
		T^1=\mathrm{d}\theta^1-\alpha\wedge\theta^2=(c_1-a)\Theta,
		\qquad
		T^2=\mathrm{d}\theta^2+\alpha\wedge\theta^1=(c_2-b)\Theta.
	\]
	因此无挠条件唯一确定 \(a=c_1,b=c_2\)。这一维数下的 \(\delta\) 既能消去所有挠率，又没有剩余的联络差；所得正是\cref{g1:prop:surface-connection-properties}的唯一度量相容无挠联络。

	但是正交标架仍可旋转。若
	\[
		g=\begin{pmatrix}\cos\varphi&-\sin\varphi\\
		\sin\varphi&\cos\varphi\end{pmatrix},
	\]
	则 \(g^{-1}J_0g=J_0\)、\(g^{-1}\mathrm{d}g=J_0\mathrm{d}\varphi\)，故 \(\alpha'=\alpha+\mathrm{d}\varphi\)。这与\cref{g1:prop:connection-frame-change}的符号一致。消去挠率并未消去旋转自由，也不能任意使 \(\alpha'=0\)：若在一个开集上能如此选择，便有 \(\mathrm{d}\alpha=0\)，由\cref{g1:prop:connection-curvature-form}必有 \(K=0\)。
\end{example}

更一般的结构可能留下非零内蕴挠率、非唯一相容联络或无法一次固定的结构群。下面只对刚才的曲面模型继续归一化：我们将使用一个由度量决定且非零的切向量来选定旋转角。这个具体构造既说明归一化怎样工作，也明确显示它在哪些点会失效。

\subsection{曲率梯度与曲面标架的归一化}
\label{g1:subsec:D03:02}

设 \((U,h)\) 为定向光滑Riemann 度量曲面的一块开区域，\(K\) 为 Gauss 曲率，并假设 \(\mathrm{d}K\) 在 \(U\) 上处处非零。这里讨论的是带度量的实二维曲面。梯度由
\[
	h(\operatorname{grad}_h K,V)=\mathrm{d}K(V)
\]
对全部切向量 \(V\) 定义；度量正定使这一向量唯一，且 \(\mathrm{d}K\ne0\) 保证它不为零。令
\begin{equation}\label{g1:eq:curvature-gradient-frame}
	q=|\operatorname{grad}_h K|_h>0,\qquad
	e_1=\frac{\operatorname{grad}_h K}{q},\qquad e_2=Je_1,
\end{equation}
其中 \(J\) 为由 \(h\) 与定向确定的正向直角旋转。于是 \(e_1\) 指向曲率增加最快的方向，\(e_2\) 的方向由定向确定；\(\operatorname{SO}(2)\) 中保持这组标架的元素只有单位元。

\ProseParagraphBreak

记对偶余标架为 \(\theta^i=h(e_i,-)\)。由梯度定义，
\[
	\mathrm{d}K=q\theta^1,\qquad \theta^1=\frac{\mathrm{d}K}{q},
	\qquad h=(\theta^1)^2+(\theta^2)^2.
\]
相应联络一形式仍取第九章的约定
\[
	\alpha(V)=h(\nabla_Ve_1,e_2),\qquad
	\nabla e_1=\alpha e_2,\qquad \nabla e_2=-\alpha e_1.
\]
因为标架已经由度量选定，\(\alpha\) 在这一步也不再含任意旋转角。

\begin{proposition}[归一化曲面标架的结构式]\label{g1:prop:normalized-surface-structure}
	设 \((U,h)\) 为定向光滑 Riemann 度量曲面，其 Gauss 曲率 \(K\) 满足 \(\mathrm dK\) 处处非零。令 \(J\) 为各切平面内的正向直角旋转，\(q=|\operatorname{grad}_h K|_h\)、\(e_1=\operatorname{grad}_h K/q\)、\(e_2=Je_1\)、\(\theta^i=h(e_i,-)\)。取 \(h\) 的 Levi--Civita 联络 \(\nabla\)，定义 \(\alpha(V)=h(\nabla_Ve_1,e_2)\)，并写 \(\alpha=a\theta^1+b\theta^2\)。则
	\begin{equation}\label{g1:eq:normalized-surface-structure}
		\begin{aligned}
			\mathrm{d}\theta^1&=\alpha\wedge\theta^2=a\theta^1\wedge\theta^2,\\
			\mathrm{d}\theta^2&=-\alpha\wedge\theta^1=b\theta^1\wedge\theta^2,\\
			\mathrm{d}\alpha&=-K\theta^1\wedge\theta^2.
		\end{aligned}
	\end{equation}
	此外，
	\begin{equation}\label{g1:eq:normalized-surface-coefficients}
		a=\frac{e_2(q)}q,\qquad
		K=e_2(a)-e_1(b)-a^2-b^2,\qquad
		[e_1,e_2]=-ae_1-be_2.
	\end{equation}
\end{proposition}
\begin{proof}
	度量相容和无挠性给出
	\[
		\begin{aligned}
			\mathrm{d}\theta^1(V,W)
			&=Vh(e_1,W)-Wh(e_1,V)-h(e_1,[V,W])\\
			&=h(\nabla_Ve_1,W)-h(\nabla_We_1,V)\\
			&=\alpha(V)\theta^2(W)-\alpha(W)\theta^2(V).
		\end{aligned}
	\]
	第二行消去了 \(h(e_1,\nabla_VW-\nabla_WV-[V,W])=0\)。所以 \(\mathrm{d}\theta^1=\alpha\wedge\theta^2\)；对 \(e_2\) 作同样计算，利用 \(\nabla e_2=-\alpha e_1\)，得到第二式。

	曲率式的符号也可直接核对。对 \(e_1\) 应用 \(R(V,W)=\nabla_V\nabla_W-\nabla_W\nabla_V-\nabla_{[V,W]}\)，沿 \(e_1\) 的乘积项抵消，余下
	\[
		R(V,W)e_1=\mathrm{d}\alpha(V,W)e_2.
	\]
	\cref{g1:prop:surface-curvature-two-dim}在正交单位基上给出 \(R(e_1,e_2)e_1=-Ke_2\)，故 \(\mathrm{d}\alpha(e_1,e_2)=-K\)，也就得到第三式。

	再对 \(\mathrm{d}K=q\theta^1\) 取外微分。因为 \(\mathrm{d}q=e_1(q)\theta^1+e_2(q)\theta^2\)，有
	\[
		0=\mathrm{d}q\wedge\theta^1+q\mathrm{d}\theta^1
		=\bigl(-e_2(q)+qa\bigr)\theta^1\wedge\theta^2,
	\]
	从而 \(a=e_2(q)/q\)。将 \(\alpha=a\theta^1+b\theta^2\) 求导，得到
	\[
		\mathrm{d}\alpha
		=\bigl(-e_2(a)+a^2+e_1(b)+b^2\bigr)\theta^1\wedge\theta^2.
	\]
	与曲率式比较便得 \(K\) 的表达。最后，\(\theta^i(e_j)\) 为常数，故 \(\mathrm{d}\theta^i(e_1,e_2)=-\theta^i([e_1,e_2])\)。前两条结构式分别给出括号的两个分量 \(-a,-b\)。
\end{proof}

这些式子中的 \(K,q,a,b\) 都由带定向的度量确定。继续沿已由度量选定的 \(e_1,e_2\) 求导，所得量也由同一几何数据决定；任意坐标偏导数则还依赖坐标选择。特别地，比较两个曲面的局部几何时，归一化余标架给出了一个明确的充分必要条件。

\begin{proposition}[归一化余标架的局部等价判据]\label{g1:prop:normalized-coframe-equivalence}
	设 \((U,h)\)、\((\widehat U,\widehat h)\) 为定向光滑Riemann 度量曲面，两个曲率函数的微分均处处非零。对第一曲面，以 \(K\) 记 Gauss 曲率，\(J\) 记切平面内的正向直角旋转，令 \(q=|\operatorname{grad}_h K|_h\)、\(e_1=\operatorname{grad}_h K/q\)、\(e_2=Je_1\)、\(\theta^i=h(e_i,-)\)；以 \(\nabla\) 记其 Levi--Civita 联络，令 \(\alpha(V)=h(\nabla_Ve_1,e_2)\)。在第二曲面上以同样的定义得到带帽记号 \(\widehat K,\widehat q,\widehat e_i,\widehat\theta^i,\widehat\alpha\)。对光滑映射 \(\Phi:U\to\widehat U\)，下列条件等价：
	\begin{equivalences}
		\item \(\Phi\) 是保持定向的局部等距。
		\item \(\Phi^*\widehat\theta^1=\theta^1\) 且 \(\Phi^*\widehat\theta^2=\theta^2\)。
	\end{equivalences}
	满足这些条件时，\(\widehat K\circ\Phi=K\)、\(\widehat q\circ\Phi=q\)，并有 \(\Phi^*\widehat\alpha=\alpha\)。
\end{proposition}
\begin{proof}
	先设 \(\Phi\) 是保持定向的局部等距。沿 \(\Phi\) 拉回目标上的度量相容无挠联络，仍满足度量相容和无挠性；\cref{g1:prop:surface-connection-properties}的唯一性说明它就是源上的联络。因此 \(\Phi\) 也保持曲率，二维曲率表达给出 \(\widehat K\circ\Phi=K\)。

	对每个切向量 \(V\)，使用梯度的定义和等距性，
	\[
		\widehat h\bigl(\Phi_*\operatorname{grad}_hK,\Phi_*V\bigr)
		=h(\operatorname{grad}_hK,V)
		=\mathrm{d}\widehat K(\Phi_*V).
	\]
	\(\Phi_*\) 在每点是同构，故 \(\Phi_*\operatorname{grad}_hK=\operatorname{grad}_{\widehat h}\widehat K\)，并保持其长度。于是 \(\Phi_*e_1=\widehat e_1\)；保持定向与度量还保证 \(\Phi_*J=\widehat J\Phi_*\)，所以 \(\Phi_*e_2=\widehat e_2\)。对偶后得到两条余标架等式。联络被保持且归一化标架对应，便有 \(\Phi^*\widehat\alpha=\alpha\)。

	反过来，两条余标架等式使 \(\mathrm{d}\Phi\) 的核为零：若 \(\Phi_*V=0\)，则 \(\theta^1(V)=\theta^2(V)=0\)，故 \(V=0\)。两侧维数同为二，所以逆函数定理使 \(\Phi\) 成为局部微分同胚。又
	\[
		\Phi^*\widehat h
		=\Phi^*\bigl((\widehat\theta^1)^2+(\widehat\theta^2)^2\bigr)
		=(\theta^1)^2+(\theta^2)^2=h,
	\]
	并且 \(\Phi^*(\widehat\theta^1\wedge\widehat\theta^2)=\theta^1\wedge\theta^2\) 保持正向面积形式。因此 \(\Phi\) 是保持定向的局部等距。
\end{proof}

判据中的等式是整个邻域上的一形式等式，不是只比较某个点的几个数。它也给出待求映射应满足的一阶方程：若 \(\theta^i=C^i_j\mathrm{d}x^j\)、\(\widehat\theta^i=\widehat C^i_a\mathrm{d}y^a\)，则
\[
	\widehat C^i_a(\Phi(x))\,\partial_j\Phi^a(x)=C^i_j(x).
\]
找到满足这些式子的局部映射，才得到等价。有限组标量相等给出必要限制，邻域上的等价还要求解这些一形式方程。

\begin{example}[一个曲率梯度标架]\label{g1:exm:normalized-warped-surface}
	在 \(x>0\) 的平面区域上取定向 \(\mathrm{d}x\wedge\mathrm{d}y>0\) 及度量
	\[
		h_f=\mathrm{d}x^2+f(x)^2\mathrm{d}y^2,\qquad f(x)>0.
	\]
	先用 \(E_1=\partial_x\)、\(E_2=f^{-1}\partial_y\) 及对偶 \(\eta^1=\mathrm{d}x\)、\(\eta^2=f\mathrm{d}y\)。由于
	\[
		\mathrm{d}\eta^1=0,\qquad
		\mathrm{d}\eta^2=\frac{f'}f\eta^1\wedge\eta^2,
	\]
	\cref{g1:exm:surface-torsion-absorption}确定联络一形式为 \(\beta=(f'/f)\eta^2=f'\mathrm{d}y\)。再算
	\[
		\mathrm{d}\beta=f''\mathrm{d}x\wedge\mathrm{d}y
		=\frac{f''}f\eta^1\wedge\eta^2,
	\]
	与曲率式比较得到 \(K=-f''/f\)。

	现在取 \(f=e^{x^2}\)。因 \(f''=(2+4x^2)f\)，有
	\[
		K=-2-4x^2,\qquad\operatorname{grad}_{h_f}K=-8x\partial_x,
		\qquad q=8x.
	\]
	所以归一化标架及余标架为
	\[
		e_1=-\partial_x,\qquad e_2=-f^{-1}\partial_y,
		\qquad \theta^1=-\mathrm{d}x,\quad\theta^2=-f\mathrm{d}y.
	\]
	这里两个向量都改变了符号，是旋转 \(\pi\)，仍保持定向。旋转角为常数，故 \(\alpha=\beta=f'\mathrm{d}y=-2x\theta^2\)，从而 \(a=0,b=-2x\)。直接核验，
	\[
		\mathrm{d}\theta^1=0,\qquad
		\mathrm{d}\theta^2=-2x\theta^1\wedge\theta^2,
		\qquad \mathrm{d}\alpha=(2+4x^2)\theta^1\wedge\theta^2.
	\]
	同时 \(e_2(q)=0\)、\(e_1(b)=2\)，故式\eqref{g1:eq:normalized-surface-coefficients}给出的曲率为 \(-2-4x^2\)，与直接计算一致。\(x=0\) 处度量仍然光滑正定，但 \(\mathrm{d}K=0\)，不能再用曲率梯度来选 \(e_1\)。
\end{example}

\begin{example}[相同的曲率和梯度长度仍不足以确定度量]\label{g1:exm:same-curvature-different-metrics}
	在同一个 \(x>0\) 区域取
	\[
		f_1(x)=e^{x^2},\qquad
		I(x)=\int_0^x e^{-2t^2}\,\mathrm{d}t,\qquad
		f_2(x)=e^{x^2}I(x).
	\]
	两函数均为正。因为 \(I'=e^{-2x^2}\)、\(I''=-4xI'\)，
	\[
		f_2''=f_1''I+2f_1'I'+f_1I''
		=(2+4x^2)f_2.
	\]
	中间两项因 \(f_1'=2xf_1\) 而相消。因此两度量 \(h_i=\mathrm{d}x^2+f_i(x)^2\mathrm{d}y^2\) 具有完全相同的函数 \(K=-2-4x^2\)，而且 \(q=8x\)。归一化标架下却有
	\[
		a_1=a_2=0,\qquad b_1=-2x,\qquad
		b_2=-2x-\frac{I'}I.
	\]
	由于 \(I'/I>0\)，两项 \(b\) 不相同。

	还可以直接排除任何局部等距，连反向等距也包括在内。若局部等距写成 \(\Phi(x,y)=(X(x,y),Y(x,y))\)，曲率必须相等。函数 \(-2-4x^2\) 在 \(x>0\) 上严格单调，所以 \(X=x\)。将此代入 \(\Phi^*h_2=h_1\)，比较 \(\mathrm{d}x^2\) 的系数得 \(1+f_2(x)^2Y_x^2=1\)，故 \(Y_x=0\)。缩小到一个开矩形后，可写 \(Y=\psi(y)\)。再比较 \(\mathrm{d}y^2\) 的系数，得到
	\[
		|\psi'(y)|=\frac{f_1(x)}{f_2(x)}=\frac1{I(x)}.
	\]
	左边只依赖 \(y\)，右边随 \(x\) 严格变化，不可能在一个开矩形上相等。因此没有这样的局部等距。相同的 \(K,q\) 留下了未被测量的横向伸缩，而 \(b\) 或完整余标架能够识别这一区别。
\end{example}

在 \(\mathrm{d}K=0\) 的点，式\eqref{g1:eq:curvature-gradient-frame}的分母为零；若 \(K\) 在开集上恒定，则整个开集都不能用这种归一化。此时可以另找非零的几何数据，或保留适当的标架自由。是否有这样的数据、剩余群是否还能约化，都须针对所选结构继续分析。

\subsection{局部等价与整体单值}
\label{g1:subsec:D03:03}

上面的标架计算比较邻域中的几何。即使每个邻域都等价于同一个齐性模型，沿环路接续这些等价时，也可能回到模型中的另一个位置或另一个标架。第十八章对平坦主联络用环路运输记录单值，\cref{g1:prop:flat-principal-monodromy}的表示取值于结构群。本节的Euclid 几何还记录位置，因此使用保向Euclid 运动群 \(\operatorname{SE}(2)=\mathbb R^2\rtimes\operatorname{SO}(2)\)，其中一个元素作用为 \(z\mapsto Qz+t\)。

\ProseParagraphBreak

下面在平圆柱、平环面与Euclid 开子域中具体构造这种接续。设 \(p:\widetilde M\to M\) 为相应模型的万有覆盖，\(\Gamma\) 为其左覆盖变换群。接续后的局部Euclid 坐标给出局部等距 \(\operatorname{dev}:\widetilde M\to\mathbb R^2\)，以及同态 \(\rho:\Gamma\to\operatorname{SE}(2)\)，满足
\begin{equation}\label{g1:eq:euclidean-developing-equivariance}
	\operatorname{dev}(\gamma\cdot\widetilde x)
	=\rho(\gamma)\operatorname{dev}(\widetilde x).
\end{equation}
这个式子比较绕路后的位置与标架变化：左边先改变覆盖片，右边在模型中施加Euclid 运动。下面在每个模型中写出 \(p,\operatorname{dev},\rho\)，从局部坐标的接续直接得到此式。更换模型中的初始Euclid 坐标 \(\operatorname{dev}'=a\circ\operatorname{dev}\)，其中 \(a\in\operatorname{SE}(2)\)，会将 \(\rho\) 改为 \(a\rho a^{-1}\)，这是将两边同时复合 \(a\) 的直接结果。

\begin{example}[平圆柱的平移单值]\label{g1:exm:flat-cylinder-developing}
	给定 \(\ell>0\)，令
	\[
		C_\ell=\mathbb R^2/\langle T\rangle,
		\qquad T(x,y)=(x+\ell,y).
	\]
	平移保持Euclid 度量，故 \(\mathrm{d}x^2+\mathrm{d}y^2\) 下降为 \(C_\ell\) 上的光滑度量。投影 \(p(x,y)=[x,y]\) 在宽度小于 \(\ell\) 的足够小开矩形上单射，给出均匀覆盖邻域；\(\mathbb R^2\) 单连通，所以它是万有覆盖。覆盖变换为 \(T^n(x,y)=(x+n\ell,y)\)，\(n\in\mathbb Z\)。取
	\[
		\operatorname{dev}(x,y)=(x,y),\qquad
		\rho(n)(u,v)=(u+n\ell,v).
	\]
	\(\rho(n+m)=\rho(n)\rho(m)\)，且 \(\operatorname{dev}\circ T^n=\rho(n)\circ\operatorname{dev}\)，故式\eqref{g1:eq:euclidean-developing-equivariance}成立。

	各小矩形都等距于平面，因而 \(K=0\)。向量场 \(\partial_x,\partial_y\) 也在平移下不变，下降成整个圆柱上的平行正交标架。因此切丛的环路平行移动全部为恒等。可是 \(\rho(1)\) 是长度为 \(\ell\) 的平移，并不平凡：只看切向量的环路平行移动，没有记录这一位置变化。局部一形式 \(\mathrm{d}x\) 同样下降为全局闭一形式 \(\xi\)，沿一周圆柱有 \(\int\xi=\ell\)，故它不是一个全局实值函数的微分。
\end{example}

\cref{g1:fig:euclidean-cylinder-strip}将一个基本带的两条竖边按相同高度对应；一次绕行在覆盖平面上留下长度为 \(\ell\) 的端点差。圆柱的周期长度可以仅用度量和环路恢复，因此不能由局部坐标选择消去。

% !TeX root = ../Book1.tex
\begin{figure}[htbp]
	\centering
	\begin{tikzpicture}[x=1.25cm,y=0.9cm,font=\small]
		\fill[GeometryBlue!7] (0,-1.5) rectangle (3,1.5);
		\draw[dashed,GeometryGray] (0,-1.5)--(3,-1.5);
		\draw[dashed,GeometryGray] (0,1.5)--(3,1.5);
		\draw[thick,GeometryBlue] (0,-1.5)--(0,1.5);
		\draw[thick,GeometryBlue] (3,-1.5)--(3,1.5);
		\draw[->,very thick,GeometryBlue] (0,-0.8)--(0,-0.15);
		\draw[->,very thick,GeometryBlue] (3,-0.8)--(3,-0.15);
		\draw[->,thick] (0,0.45)--(3,0.45);
		\fill (0,0.45) circle (1.3pt);
		\fill (3,0.45) circle (1.3pt);
		\node[above] at (1.5,0.45) {\(\ell\)};
		\node[left] at (0,0.45) {\((0,t)\)};
		\node[right] at (3,0.45) {\((\ell,t)\)};
		\draw[->,GeometryGray] (0,1.85)--(3,1.85);
		\node[above] at (1.5,1.85) {\(T(s,t)=(s+\ell,t)\)};
		\node[below] at (0,-1.5) {\(s=0\)};
		\node[below] at (3,-1.5) {\(s=\ell\)};
	\end{tikzpicture}
	\caption[平圆柱的基本带与平移周期]{平圆柱由基本带的两条竖边按同高点粘合得到，两边箭头同向。横线段投影为一次闭路，在覆盖平面上的端点相差平移 \(\ell\)。上下虚线只表示绘图截断，实际基本带沿 \(t\) 方向无限延伸。}
	\label{g1:fig:euclidean-cylinder-strip}
\end{figure}


\begin{proposition}[平圆柱的最短非平凡环路]\label{g1:prop:cylinder-period-length}
	设 \(\ell>0\)，以平移 \((x,y)\mapsto(x+\ell,y)\) 生成的群作用定义圆柱 \(C_\ell=\mathbb R^2/(\ell\mathbb Z\times\{0\})\)，并赋予标准平度量 \(\mathrm dx^2+\mathrm dy^2\) 下降的商度量。在 \(C_\ell\) 上，所有不可缩的分段 \(C^1\) 环路的长度至少为 \(\ell\)，且存在长度恰为 \(\ell\) 的不可缩环路。对任意 \(\ell'>0\)，以同样方式定义带商度量的 \(C_{\ell'}\)。因此若 \(C_\ell\) 与 \(C_{\ell'}\) 整体等距，则 \(\ell=\ell'\)。
\end{proposition}
\begin{proof}
	将一条环路从 \((x_0,y_0)\) 提升到 \(\mathbb R^2\)，终点必为 \((x_0+n\ell,y_0)\)，其中 \(n\in\mathbb Z\)。由\cref{g1:cor:closed-lift-criterion}，万有覆盖上的提升闭合当且仅当原环路可缩；这里 \(\mathbb R^2\) 的基本群平凡，所以不可缩恰等价于 \(n\ne0\)。投影是局部等距，故原曲线与提升曲线长度相同。提升曲线的Euclid 长度不小于端点距离，因而
	\[
		L(\gamma)\ge |n|\ell\ge\ell.
	\]
	水平线段 \(t\mapsto(x_0+t\ell,y_0)\)、\(0\le t\le1\)，投影后是一条环路，提升的端点差对应 \(n=1\)，所以它不可缩且长度恰为 \(\ell\)。整体等距是微分同胚，保持环路的可缩性与曲线长度，故必须保持这个最小值。
\end{proof}

所以所有 \(C_\ell\) 的局部曲率都相同，切向量的平行移动单值也都平凡，仍有连续变化的整体周期数据。环面上的现象相似：若 \(\Lambda\subset\mathbb R^2\) 是满格，则平环面 \(\mathbb R^2/\Lambda\) 的覆盖空间仍为 \(\mathbb R^2\)，可取 \(\operatorname{dev}(z)=z\)、\(\rho(\lambda)(z)=z+\lambda\)。格向量记录两个方向的平移周期，不能从 \(K=0\) 中恢复。

\ProseParagraphBreak

即使完整Euclid 单值 \(\rho\) 平凡，也只说明展开映射在各覆盖片上的取值相同。由式\eqref{g1:eq:euclidean-developing-equivariance}，此时可定义
\[
	\overline{\operatorname{dev}}(p(\widetilde x))
	=\operatorname{dev}(\widetilde x).
\]
同一纤维的点相差覆盖变换，故定义无关；在均匀覆盖邻域上使用一个局部逆即可看出它光滑且为局部等距。这个结论没有保证 \(\overline{\operatorname{dev}}\) 单射，也没有保证它覆盖整个 \(\mathbb R^2\)。

\begin{example}[平凡单值不保证整体等价于平面]\label{g1:exm:trivial-monodromy-incomplete-model}
	单位开圆盘 \(B\subset\mathbb R^2\) 带Euclid 度量时，万有覆盖就是自身，\(\operatorname{dev}:B\hookrightarrow\mathbb R^2\) 为包含映射，\(\rho\) 平凡。它是单值且单射的局部等距，但不满射。圆盘凸，任意两点之间的Euclid 线段仍在圆盘内，所以内在距离等于Euclid 距离，直径为 \(2\)；整个Euclid 平面的直径无界。因此圆盘也不可能通过另一个映射整体等距于平面。

	单射性也可能失败。取 \(N=(0,\infty)\times\mathbb R\)，赋予度量 \(\mathrm{d}r^2+r^2\mathrm{d}t^2\)，并令
	\[
		D(r,t)=(r\cos t,r\sin t).
	\]
	其两个偏导线性无关，且 \(D^*(\mathrm{d}x^2+\mathrm{d}y^2)=\mathrm{d}r^2+r^2\mathrm{d}t^2\)，所以 \(D\) 是局部等距。\(N\) 是凸开集，故单连通；把它作为底流形时，万有覆盖可取恒等映射，单值仍平凡。然而 \(D(r,t+2\pi)=D(r,t)\)，所以展开映射并不单射，像也只是穿孔平面。
\end{example}

在这些例子中，局部结构方程确定邻域中的模型，单值确定绕行后的识别。当单值平凡、展开映射下降到底流形后，所得局部等距若再是到整个模型的双射，便给出整体同构。因而还需分别考察它的单射性与像。

\begin{chaptersummary}
	一个 \(G\)-结构规定了每点允许的标架，典范形式把标架坐标同真实切向量联系起来。相容联络的差是张量，它通过交替化改变挠率；在相应商空间中留下的类就是内蕴挠率。正交结构的内蕴挠率为零且无挠相容联络唯一，但其曲率可以非零。几乎复和接触模型则显示，适应标架的括号可能产生不能由联络选择消去的分量。

	\ProseParagraphBreak

	Cartan 联络在主 \(H\) 丛的每个切空间上给出到模型Lie 代数 \(\mathfrak l\) 的同构，既再现竖直方向，又记录底空间方向；它的核为零。Ehresmann 主联络则用 \(\mathfrak h\) 值形式的核选择水平空间。只有在给出相容的分解后，才能把 Cartan 形式拆成主联络部分与切向部分。曲率为零时，结构方程可以局部积分成齐性模型，路径展开进一步描述了这种比较怎样沿路延伸。

	\ProseParagraphBreak

	等价问题需要同时处理标架自由和结构函数。在曲面上，非零的曲率梯度能够选出一组规范的正交标架，随后由联络系数及其变化检验度量；曲率函数相同本身仍不足够。这种归一化有明确的适用区域，不能直接跨过梯度为零的点。平圆柱的周期及模型中的开子域则表明，局部模型、整体单值和全局同构是不同层次的结论。
\end{chaptersummary}

\BookExercises

\begin{exercise}
	\label{g1:ex:intrinsic-torsion-global-frame}
	在 \(\mathbb R^2\) 上取标架 \(e_1=\partial_x\)、\(e_2=e^x\partial_y\)。将这个指定标架视为 \(G=\{I\}\) 的结构群约化。
	\begin{enumerate}
		\item 求对偶余标架以及唯一适应联络的挠率。
		\item 证明不存在任何非空开集上的坐标 \((u,v)\)，使同时有 \(e_1=\partial_u\)、\(e_2=\partial_v\)。
		\item 将第二个标架向量改为 \(\widetilde e_2=e^{-x}e_2\)，所得标架的内蕴挠率为零。这是否与内蕴挠率不随适应标架改变相矛盾？
	\end{enumerate}
\end{exercise}
\begin{solution}
	(1) 对偶余标架为 \(\vartheta^1=\mathrm{d}x\)、\(\vartheta^2=e^{-x}\mathrm{d}y\)。由于 \(\mathfrak g=0\)，唯一适应联络满足 \(\nabla e_1=\nabla e_2=0\)。其挠率分量为
	\[
		\mathrm{d}\vartheta^1=0,\qquad
		\mathrm{d}\vartheta^2=-e^{-x}\mathrm{d}x\wedge\mathrm{d}y
		=-\vartheta^1\wedge\vartheta^2.
	\]
	因此 \(T(e_1,e_2)=-e_2\)。这里 \(C=V^*\otimes\mathfrak g=0\)，\(\operatorname{im}\delta=0\)，所以这一挠率本身就是内蕴挠率，处处非零。

	(2) 直接计算 \([e_1,e_2]=e^x\partial_y=e_2\)，处处不为零。而同一坐标系中的两个坐标向量场必定对易，因此不可能在任何非空开集同时满足题设两式。

	(3) 新标架为 \((\partial_x,\partial_y)\)，对偶形式均闭，其唯一适应联络无挠。但从旧标架到新标架的矩阵为 \(\operatorname{diag}(1,e^{-x})\)，并非恒等群 \(\{I\}\) 值函数，因此不是同一约化内允许的适应标架变换。题目改变了 \(G\)-结构本身，内蕴挠率可以随之改变。
\end{solution}

\begin{exercise}
	\label{g1:ex:conformal-coframe-torsion-correction}
	设 \(U\subset\mathbb R^2\) 为开集，\(f\in C^\infty(U)\)，在 \(U\) 上取度量
	\[
		g=e^{2f}(\mathrm{d}x^2+\mathrm{d}y^2)
	\]
	及正交单位余标架 \(\vartheta^1=e^f\mathrm{d}x\)、\(\vartheta^2=e^f\mathrm{d}y\)。在其对偶标架中，把度量联络矩阵写为
	\[
		a=\begin{pmatrix}0&-\eta\\\eta&0\end{pmatrix}.
	\]
	从无挠结构方程求出一形式 \(\eta\)，并在 \(f(x,y)=x^2+y\) 时写出结果。说明无挠为何不要求这两个余标架分量分别闭。
\end{exercise}
\begin{solution}
	写 \(\eta=p\,\mathrm{d}x+q\,\mathrm{d}y\)。外微分为
	\[
		\mathrm{d}\vartheta^1=-e^ff_y\,\mathrm{d}x\wedge\mathrm{d}y,
		\qquad
		\mathrm{d}\vartheta^2=e^ff_x\,\mathrm{d}x\wedge\mathrm{d}y.
	\]
	方程 \(\mathrm{d}\vartheta+a\wedge\vartheta=0\) 的两个分量分别为
	\[
		\mathrm{d}\vartheta^1-\eta\wedge\vartheta^2
		=-e^f(f_y+p)\,\mathrm{d}x\wedge\mathrm{d}y=0,
	\]
	\[
		\mathrm{d}\vartheta^2+\eta\wedge\vartheta^1
		=e^f(f_x-q)\,\mathrm{d}x\wedge\mathrm{d}y=0.
	\]
	故唯一解为 \(p=-f_y\)、\(q=f_x\)，即
	\[
		\eta=-f_y\,\mathrm{d}x+f_x\,\mathrm{d}y.
	\]
	当 \(f=x^2+y\) 时，\(\eta=-\mathrm{d}x+2x\,\mathrm{d}y\)。此时 \(\mathrm{d}\vartheta^1\) 处处非零，但它恰被 \(-\eta\wedge\vartheta^2\) 抵消。无挠要求 \(\mathrm{d}\vartheta+a\wedge\vartheta=0\)，所以余标架变化与联络作用要合在一起；只有联络矩阵也为零时，无挠才要求各余标架分量闭。
\end{solution}

\begin{exercise}
	\label{g1:ex:almost-complex-variable-obstruction}
	设 \(I\subset\mathbb R\) 是非空开区间，\(f\in C^\infty(I)\)。在 \(I\times\mathbb R^3\) 上取
	\[
		e_1=\partial_x,\quad e_2=\partial_y,\quad
		e_3=\partial_z,\quad e_4=\partial_t+f(x)\partial_y,
	\]
	并按 \(Je_1=e_2\)、\(Je_3=e_4\)、\(J^2=-I\) 定义几乎复结构。
	\begin{enumerate}
		\item 计算 \(\mathcal K_J(e_1-\mathrm i e_2,e_3-\mathrm i e_4)\)。
		\item 证明这个结构在整个 \(I\times\mathbb R^3\) 上有无挠适应联络，当且仅当 \(f\) 为常数。
		\item 当 \(f=c\) 为常数时，给出实坐标，使 \(J\) 在坐标标架中的矩阵为标准常矩阵。
	\end{enumerate}
\end{exercise}
\begin{solution}
	(1) 唯一可能非零的基本括号为 \([e_1,e_4]=f'(x)e_2\)。令 \(Z=e_1-\mathrm i e_2\)、\(W=e_3-\mathrm i e_4\)，则
	\[
		[Z,W]=-\mathrm i f'(x)e_2
		=\frac{f'(x)}2(Z-\overline Z),
	\]
	故 \(\mathcal K_J(Z,W)=-\tfrac12 f'(x)\overline Z\)。

	(2) 若有无挠适应联络，由\cref{g1:prop:almost-complex-intrinsic-component}，上一式处处为零，所以在 \(I\) 上有 \(f'=0\)。区间连通，故 \(f\) 为常数。反之，若 \(f=c\)，所有基本括号都为零。规定 \(\nabla e_i=0\)，得到保持 \(J\) 的联络，且 \(T(e_i,e_j)=-[e_i,e_j]=0\)。由于挠率是张量，这就说明其在所有输入上为零。

	(3) 取新坐标 \((u,v,w,s)=(x,y-ct,z,t)\)，逆变换为 \((x,y,z,t)=(u,v+cs,w,s)\)。因此
	\[
		\partial_u=e_1,\quad \partial_v=e_2,\quad
		\partial_w=e_3,\quad \partial_s=e_4.
	\]
	在此坐标中，\(J\partial_u=\partial_v\)、\(J\partial_w=\partial_s\)，故确为标准常矩阵。这个结论来自本模型的显式换坐标，不需要调用一般的几乎复可积性定理。
\end{solution}

\begin{exercise}
	\label{g1:ex:contact-form-scaling-intrinsic}
	在 \(\mathbb R^3\) 上令 \(\alpha=\mathrm{d}z-x\,\mathrm{d}y\)、\(D=\ker\alpha\)，并令 \(h\) 为处处非零光滑实函数，\(\widetilde\alpha=h\alpha\)。
	\begin{enumerate}
		\item 证明两一形式定义同一个 \(G_D\)-结构，并计算 \(\widetilde\alpha\wedge\mathrm{d}\widetilde\alpha\)。
		\item 对 \(X,Y\in\Gamma(D)\)，证明
		\[
			\alpha([X,Y])=-\mathrm{d}\alpha(X,Y),\qquad
			\widetilde\alpha([X,Y])=h\alpha([X,Y]).
		\]
		解释为什么右边的倍数变化没有改变商丛值括号张量 \(\mathcal B\)。
		\item 设 \((A,B,C)\) 为正文中的标架。求将它改成满足 \(u^*\widetilde\alpha=\varepsilon^3\) 的标架的一种方法，并说明该变换是否保持严格条件 \(u^*\alpha=\varepsilon^3\)。
	\end{enumerate}
\end{exercise}
\begin{solution}
	(1) \(h\) 处处非零，故 \(\ker\widetilde\alpha=\ker\alpha=D\)。\(P_D\) 只由标架的前两列是否张成 \(D\) 决定，所以约化相同。又
	\[
		\mathrm{d}\widetilde\alpha=\mathrm{d}h\wedge\alpha+h\,\mathrm{d}\alpha,
	\]
	从而
	\[
		\widetilde\alpha\wedge\mathrm{d}\widetilde\alpha
		=h^2\alpha\wedge\mathrm{d}\alpha
		=-h^2\mathrm{d}x\wedge\mathrm{d}y\wedge\mathrm{d}z\neq0.
	\]
	其中含两次 \(\alpha\) 的项为零。

	(2) 由一形式外微分公式及 \(\alpha(X)=\alpha(Y)=0\)，有 \(\mathrm{d}\alpha(X,Y)=-\alpha([X,Y])\)。第二式直接由 \(\widetilde\alpha=h\alpha\) 得出。这里 \(\mathcal B(X,Y)\) 是 \([X,Y]\) 在 \(\mathrm{T}M/D\) 中的类，从未用 \(\alpha\) 选择数值坐标。\(\alpha\) 诱导同构 \(\mathrm{T}M/D\to\mathbb R\)，\(\widetilde\alpha\) 则把这个同构乘以 \(h\)。变的是同一个商丛元素的标量表示，商丛值张量本身不变。

	(3) 取 \((A,B,C/h)\)。前两列仍属于 \(D\)，且 \(\widetilde\alpha(C/h)=1\)，所以满足新形式的严格条件。相对于原标架，变换矩阵为 \(\operatorname{diag}(1,1,h^{-1})\)，属于 \(G_D\)。但 \(\alpha(C/h)=h^{-1}\)，故它只有在 \(h=1\) 的点才保持原来的严格条件；要在整个区域保持，则必须有 \(h\equiv1\)。这正体现了保留分布与保留一个指定一形式的区别。
\end{solution}

\begin{exercise}\label{g1:ex:cartan-trivial-bundle}
	设 \(L\) 为有限维实Lie 群，\(H\subseteq L\) 为闭子群，\(U\) 为无边界光滑流形。给定 \(B\in\Omega^1(U;\mathfrak l)\)，在主右 \(H\) 丛 \(P=U\times H\) 上定义
	\[
		\omega_{(x,h)}(v,\eta)
		=\operatorname{Ad}_{h^{-1}}B_x(v)+(\theta_H)_h(\eta).
	\]
	证明：\(\omega\) 为模型 \(L/H\) 的 Cartan 联络，当且仅当每个 \(x\in U\) 处
	\[
		B_x\bmod\mathfrak h:\mathrm{T}_xU\longrightarrow\mathfrak l/\mathfrak h
	\]
	为同构。并证明，在这一条件下，\(\omega\) 平坦当且仅当
	\[
		\mathrm{d}B+\tfrac12[B\wedge B]=0.
	\]
	这是否要求 \(B_x:\mathrm{T}_xU\to\mathfrak l\) 本身为同构？
\end{exercise}
\begin{solution}
	固定 \(k\in H\)。右作用为 \((x,h)k=(x,hk)\)，Maurer--Cartan 形式的右平移律给出
	\[
		\begin{aligned}
			(R_k^*\omega)_{(x,h)}(v,\eta)
			&=\operatorname{Ad}_{(hk)^{-1}}B_x(v)
			+(\theta_H)_{hk}((\mathrm{d}R_k)_h\eta)\\
			&=\operatorname{Ad}_{k^{-1}}
			\bigl(\operatorname{Ad}_{h^{-1}}B_x(v)+(\theta_H)_h(\eta)\bigr).
		\end{aligned}
	\]
	故等变性对任意 \(B\) 成立。基本向量为 \((0,X_h^L)\)，代入即得 \(\omega(0,X_h^L)=X\)，故垂直再现也总成立。

	先假定商映射为同构。若 \(\omega(v,\eta)=0\)，取模 \(\mathfrak h\)，得
	\[
		\operatorname{Ad}_{h^{-1}}(B_x(v)+\mathfrak h)=0,
	\]
	从而 \(v=0\)。随后 \(\theta_H(\eta)=0\) 给出 \(\eta=0\)。商同构还给出 \(\dim U+\dim H=\dim L\)，所以 \(\omega\) 为同维空间之间的单射，因而是同构。

	反之，若 \(\omega\) 为 Cartan 联络，在 \((x,e)\) 处它是
	\[
		(v,X)\longmapsto B_x(v)+X
	\]
	这个线性同构。任意 \(Y\in\mathfrak l\) 可唯一如此表示，取商说明 \(B_x\bmod\mathfrak h\) 满射。若 \(B_x(v)\in\mathfrak h\)，则 \((v,-B_x(v))\) 被映到零，同构性使 \(v=0\)，故商映射也单射。

	令 \(s(x)=(x,e)\)，则 \(s^*\omega=B\)，从而
	\[
		s^*\Omega=\mathrm{d}B+\tfrac12[B\wedge B].
	\]
	若 \(\Omega=0\)，右端当然为零。反过来，若 \(s^*\Omega=0\)，则两个输入都沿截面切向时，曲率的取值为零；水平性使插入任何垂直分量的项也为零。由于截面方向与垂直方向张成 \(\mathrm{T}_{(x,e)}P\)，曲率在 \(s(U)\) 上为零；右等变性再将零传到 \((x,h)\) 的每一点。

	\(B_x\) 本身不必、通常也不可能满射到 \(\mathfrak l\)：若 \(\dim H>0\)，源维数比目标小 \(\dim H\)。它只需在取商后为同构；缺少的 \(\mathfrak h\) 方向由 \(\theta_H\) 提供。
\end{solution}

\begin{exercise}\label{g1:ex:cartan-spherical-model}
	沿用\cref{g1:exm:cartan-euclidean-surface}的定向二维Riemann 流形及正向正交标架丛 \(P\)，保留 \(\alpha,\theta^1,\theta^2,J_0\) 的约定。现在取
	\[
		L=\operatorname{SO}(3),\qquad
		H=\{\operatorname{diag}(h,1):h\in\operatorname{SO}(2)\},
	\]
	并在同一个 \(P\) 上定义 \(\mathfrak{so}(3)\) 值形式
	\[
		\widehat\omega=
		\begin{pmatrix}
			0&-\alpha&\theta^1\\
			\alpha&0&\theta^2\\
			-\theta^1&-\theta^2&0
		\end{pmatrix}.
	\]
	证明它是模型为单位球面 \(\operatorname{SO}(3)/\operatorname{SO}(2)\) 的 Cartan 联络，直接计算其曲率，并判定它何时平坦。解释为什么单位球面的这个 Cartan 几何平坦，与该球面的Euclid  Cartan 几何不平坦并不矛盾。
\end{exercise}
\begin{solution}
	先将 \(P\) 的结构群通过 \(h\mapsto\operatorname{diag}(h,1)\) 识别为给定的 \(H\)。记 \(A=J_0\alpha\) 和 \(\theta=(\theta^1,\theta^2)^{\mathsf T}\)，则
	\[
		\widehat\omega=
		\begin{pmatrix}A&\theta\\-\theta^{\mathsf T}&0\end{pmatrix}.
	\]
	\(\widehat\omega\) 反对称，因而取值于 \(\mathfrak{so}(3)\)。主联络满足 \(R_h^*A=h^{-1}Ah\)，焊接形式满足 \(R_h^*\theta=h^{-1}\theta\)。由于 \(h^{-1}=h^{\mathsf T}\)，整个分块矩阵满足
	\[
		R_h^*\widehat\omega
		=\operatorname{diag}(h,1)^{-1}\widehat\omega
		\operatorname{diag}(h,1).
	\]
	垂直输入时 \(\theta=0\)，而 \(A\) 再现 \(\mathfrak{so}(2)\)，故再现公理成立。若 \(\widehat\omega(V)=0\)，先由 \(\theta(V)=0\) 知 \(V\) 垂直，再由 \(A(V)=0\) 得 \(V=0\)。\(\dim P=\dim\mathfrak{so}(3)=3\)，所以逐点为同构。

	\(\operatorname{SO}(3)\) 在单位球面上传递，固定北极的稳定子恰为上述 \(H\)。由\cref{g1:prop:transitive-action-model}，商模型是单位球面。矩阵Lie 代数中的曲率为
	\(\widehat\Omega=\mathrm{d}\widehat\omega+\widehat\omega\wedge\widehat\omega\)。逐项相乘得
	\[
		\widehat\Omega=
		\begin{pmatrix}
			0&-\mathrm{d}\alpha-\theta^1\wedge\theta^2&
			\mathrm{d}\theta^1-\alpha\wedge\theta^2\\
			\mathrm{d}\alpha+\theta^1\wedge\theta^2&0&
			\mathrm{d}\theta^2+\alpha\wedge\theta^1\\
			-\mathrm{d}\theta^1+\alpha\wedge\theta^2&
			-\mathrm{d}\theta^2-\alpha\wedge\theta^1&0
		\end{pmatrix}.
	\]
	无挠结构方程使最后一列和最后一行的非对角项为零。又有 \(\mathrm{d}\alpha=-(K\circ\pi)\theta^1\wedge\theta^2\)，所以
	\[
		\widehat\Omega=
		\begin{pmatrix}
			0&(K\circ\pi-1)\theta^1\wedge\theta^2&0\\
			(1-K\circ\pi)\theta^1\wedge\theta^2&0&0\\
			0&0&0
		\end{pmatrix}.
	\]
	焊接形式在底空间方向给出余标架，故 \(\theta^1\wedge\theta^2\) 在这些方向上非零。因此球面模型的 Cartan 曲率为零，当且仅当 \(K\equiv1\)。

	同一底空间和同一主 \(\operatorname{SO}(2)\) 丛可以承载不同模型的 Cartan 联络。Euclid 模型中两个平移方向的括号为零，曲率的旋转部分是 \(J_0\,\mathrm{d}\alpha\)；球面模型中对应方向的括号产生附加项 \(J_0\theta^1\wedge\theta^2\)。单位球面与球面模型局部相同，所以后者平坦；它与Euclid 平面局部不等距，所以前者不平坦。平坦性所比较的模型已经改变。
\end{solution}

\begin{exercise}\label{g1:ex:cartan-development-path-dependence}
	取加法Lie 群 \(L=\mathbb R^2\)、平凡子群 \(H=\{0\}\)，于是主丛总空间就是 \(P=M=\mathbb R^2\)。定义
	\[
		\omega=\begin{pmatrix}\mathrm{d}x\\e^x\,\mathrm{d}y\end{pmatrix}.
	\]
	证明这给出一个非平坦 Cartan 几何。比较从 \((0,0)\) 到 \((1,1)\) 的两条折线路径：第一条先沿 \(x\) 轴到 \((1,0)\)，再沿竖直方向到达终点；第二条先沿 \(y\) 轴到 \((0,1)\)，再沿水平方向到达终点。求各自展开路径的终点，说明在单连通底空间上，曲率不为零时为何仍不能只由路径终点定义展开映射。
\end{exercise}
\begin{solution}
	在坐标基下，\(\omega\) 的矩阵为 \(\operatorname{diag}(1,e^x)\)，行列式始终为正，所以逐点可逆。结构群平凡，等变性和再现条件没有额外要求，故 \(\omega\) 是 Cartan 联络。加法Lie 代数的括号为零，因而
	\[
		\Omega=\mathrm{d}\omega
		=\begin{pmatrix}0\\e^x\,\mathrm{d}x\wedge\mathrm{d}y\end{pmatrix}\ne0.
	\]
	模型上的 Maurer--Cartan 形式是恒等余标架。展开 \(g=(u,v)\) 满足
	\[
		u'=x',\qquad v'=e^{x(t)}y',\qquad
		(u(0),v(0))=(0,0).
	\]
	故 \(u(1)=1\)，而 \(v(1)=\int_\gamma e^x\,\mathrm{d}y\)。

	沿第一条路径，水平段贡献为零，竖直段 \(x=1\)，所以 \(v(1)=e\)，展开终点是 \((1,e)\)。沿第二条路径，竖直段 \(x=0\)，随后水平段没有贡献，所以展开终点是 \((1,1)\)。两条路径在 \(\mathbb R^2\) 中定端点同伦，却有不同的展开终点；单连通性不能消除曲率造成的路径依赖。

	差值还可以直接由曲率看出。将第一条路径与第二条的反向拼接，得到单位正方形的正向边界；Stokes 公式给出
	\[
		e-1
		=\int_{\partial[0,1]^2}e^x\,\mathrm{d}y
		=\int_{[0,1]^2}e^x\,\mathrm{d}x\wedge\mathrm{d}y.
	\]
	在这个交换模型中，路径展开的差恰好由曲率的面积积分表示。
\end{solution}

\begin{exercise}\label{g1:ex:cartan-one-dimensional-completeness}
	取 \(L=\mathbb R\)、\(H=\{0\}\)，并在 \(M=\mathbb R\) 上给定 \(\omega=f(x)\,\mathrm{d}x\)，其中 \(f:\mathbb R\to(0,\infty)\) 光滑。证明这总是平坦 Cartan 联络。令
	\[
		F(x)=\int_0^x f(s)\,\mathrm{d}s.
	\]
	求从 \(x_0\) 出发的路径展开，刻画 \(F\) 的像，并证明由 \(\omega(E)=1\) 定义的向量场 \(E\) 完备，当且仅当
	\[
		\int_0^\infty f(s)\,\mathrm{d}s=\infty,
		\qquad
		\int_{-\infty}^0 f(s)\,\mathrm{d}s=\infty.
	\]
	分别讨论 \(f=1\) 与 \(f=e^x\)，说明即使底空间是单连通的整条实直线，平坦性也不保证所有模型运动都能无限延拓。
\end{exercise}
\begin{solution}
	\(f>0\) 保证 \(\omega\) 逐点可逆，另外两条 Cartan 公理在平凡结构群下自动成立。加法Lie 代数的括号为零，而一维流形上二形式为零，故 \(\Omega=\mathrm{d}(f\,\mathrm{d}x)=0\)。

	若 \(\gamma(t)=x(t)\) 且 \(x(0)=x_0\)，展开方程为 \(g'=f(x)x'\)、\(g(0)=0\)。由链式法则，
	\[
		g(t)=F(x(t))-F(x_0).
	\]
	因此此处展开确实只依赖当前终点。又因为 \(F'=f>0\)，\(F\) 严格递增且处处为局部微分同胚。记扩展实数
	\[
		a=-\int_{-\infty}^0 f(s)\,\mathrm{d}s,\qquad
		b=\int_0^\infty f(s)\,\mathrm{d}s.
	\]
	单调性及积分定义给出 \(\lim_{x\to-\infty}F(x)=a\)、\(\lim_{x\to+\infty}F(x)=b\)。由介值定理，\(F\) 将 \(\mathbb R\) 双射到 \((a,b)\)；局部光滑逆拼接成光滑全局逆，故它是微分同胚。

	由 \(\omega(E)=1\) 得 \(E=f^{-1}\partial_x\)。其过 \(x_0\) 的积分曲线满足
	\[
		F(x(t))=F(x_0)+t,\qquad
		x(t)=F^{-1}(F(x_0)+t).
	\]
	该式恰在
	\[
		a-F(x_0)<t<b-F(x_0)
	\]
	上有定义。若积分曲线能在 \(\mathbb R\) 中延拓到某个有限端点，连续性会使该端点的 \(F\) 值属于 \((a,b)\)，与它等于 \(a\) 或 \(b\) 矛盾，所以这是最大定义区间。于是 \(E\) 的所有积分曲线在整个实时间轴上存在，当且仅当 \(a=-\infty\)、\(b=\infty\)，即题述两个积分均发散。

	当 \(f=1\) 时，\(F(x)=x\)，\(E=\partial_x\) 完备。当 \(f=e^x\) 时，
	\[
		F(x)=e^x-1,\qquad
		F(\mathbb R)=(-1,\infty),\qquad
		x(t)=\log(e^{x_0}+t).
	\]
	这条积分曲线只在 \(t>-e^{x_0}\) 上存在，故 \(E=e^{-x}\partial_x\) 不完备。底空间的拓扑与平坦性均没有改变，但模型中的匀速负向运动在有限时间后离开了 \(F(\mathbb R)\)。
\end{solution}

\begin{exercise}\label{g1:ex:rotating-plane-frame-gauge}
	在定向Euclid 平面上从标架 \((\partial_x,\partial_y)\) 出发，将标架在点 \((x,y)\) 处旋转角度 \(x\)。求新的余标架 \(\theta^1,\theta^2\) 和联络一形式 \(\alpha\)，并直接核验
	\[
		\mathrm{d}\theta^1=\alpha\wedge\theta^2,\qquad
		\mathrm{d}\theta^2=-\alpha\wedge\theta^1,\qquad
		\mathrm{d}\alpha=0.
	\]
	说明非零的 \(\alpha\) 为何不表示度量发生改变，也不表示曲率非零。
\end{exercise}
\begin{solution}
	旋转后的标架为
	\[
		e_1=\cos x\,\partial_x+\sin x\,\partial_y,\qquad
		e_2=-\sin x\,\partial_x+\cos x\,\partial_y.
	\]
	对偶余标架按旋转矩阵的逆变换，得到
	\[
		\theta^1=\cos x\,\mathrm{d}x+\sin x\,\mathrm{d}y,\qquad
		\theta^2=-\sin x\,\mathrm{d}x+\cos x\,\mathrm{d}y.
	\]
	原标架的联络一形式为零，换架公式给出 \(\alpha=\mathrm{d}x\)。也可直接计算 \(\nabla_Ve_1=\mathrm{d}x(V)e_2\) 得到相同结果。求外微分，
	\[
		\begin{aligned}
			\mathrm{d}\theta^1&=\cos x\,\mathrm{d}x\wedge\mathrm{d}y
			=\mathrm{d}x\wedge\theta^2,\\
			\mathrm{d}\theta^2&=-\sin x\,\mathrm{d}x\wedge\mathrm{d}y
			=-\mathrm{d}x\wedge\theta^1,
		\end{aligned}
	\]
	而 \(\mathrm{d}\alpha=\mathrm{d}(\mathrm{d}x)=0\)。因此挠率与曲率仍为零。

	将两个余标架平方相加，混合项抵消，得到 \((\theta^1)^2+(\theta^2)^2=\mathrm{d}x^2+\mathrm{d}y^2\)。度量没有改变；非零的 \(\alpha\) 记录所选标架随位置旋转的速率。曲率比较不同方向上的变化能否相容，这里由 \(\mathrm{d}\alpha=0\) 表示，不能仅由 \(\alpha\) 是否为零判断。
\end{solution}

\begin{exercise}\label{g1:ex:normalized-curvature-hessian}
	设 \((U,h)\) 为定向光滑Riemann 度量曲面，\(\mathrm{d}K\) 处处非零。使用式\eqref{g1:eq:curvature-gradient-frame}的归一化标架，并写 \(\alpha=a\theta^1+b\theta^2\)。对
	\[
		\operatorname{Hess}K(V,W)=h(\nabla_V\operatorname{grad}_hK,W),
	\]
	证明其在 \(e_1,e_2\) 下的矩阵为
	\[
		\begin{pmatrix}e_1(q)&e_2(q)\\e_2(q)&qb\end{pmatrix}.
	\]
	再证明曲率水平集沿单位切向量 \(e_2\) 定向时，其有向测地曲率等于 \(b\)。这里采用第九章的约定 \(k_g=h(\nabla_TT,JT)\)。
\end{exercise}
\begin{solution}
	由 \(\operatorname{grad}_hK=qe_1\) 与联络公式，
	\[
		\nabla_{e_1}\operatorname{grad}_hK=e_1(q)e_1+qa e_2,
		\qquad
		\nabla_{e_2}\operatorname{grad}_hK=e_2(q)e_1+qb e_2.
	\]
	分别与 \(e_1,e_2\) 作内积，先得到矩阵
	\[
		\begin{pmatrix}e_1(q)&qa\\e_2(q)&qb\end{pmatrix}.
	\]
	由式\eqref{g1:eq:normalized-surface-coefficients}中的 \(qa=e_2(q)\)，即得所求对称矩阵。也可以由无挠性直接检查对称性：
	\[
		\operatorname{Hess}K(V,W)-\operatorname{Hess}K(W,V)
		=VW(K)-WV(K)-[V,W](K)=0.
	\]
	这里展开 Hessian 时用了 \(\operatorname{Hess}K(V,W)=V(WK)-(\nabla_VW)K\)。

	因为 \(\mathrm{d}K\ne0\)，每个曲率水平集在局部是一条光滑曲线；\(e_2(K)=0\) 且 \(|e_2|_h=1\)，故可按 \(e_2\) 定向并用弧长参数化。其切向量 \(T=e_2\) 满足
	\[
		\nabla_TT=\nabla_{e_2}e_2=-b e_1,\qquad JT=Je_2=-e_1.
	\]
	于是 \(k_g=h(-be_1,-e_1)=b\)。因此归一化后的系数 \(b\) 同时记录 Hessian 的横向分量和曲率水平线的有向测地弯曲。
\end{solution}

\begin{exercise}\label{g1:ex:warped-surface-local-symmetries}
	在半平面 \(x>0\) 上赋予度量 \(h=\mathrm{d}x^2+e^{2x^2}\mathrm{d}y^2\) 及定向 \(\mathrm{d}x\wedge\mathrm{d}y>0\)。证明任意保持定向的局部自等距，在每个足够小的连通邻域内都形如
	\[
		\Phi(x,y)=(x,y+c)
	\]
	其中 \(c\) 为常数；反之，这些映射确实是等距。再求反转定向的局部自等距。
\end{exercise}
\begin{solution}
	\cref{g1:exm:normalized-warped-surface}已计算 \(K=-2-4x^2\)，在 \(x>0\) 上严格单调。若 \(\Phi=(X,Y)\) 是局部自等距，曲率保持性给出 \(-2-4X^2=-2-4x^2\)。由于源、目标的第一坐标都为正，只能有 \(X=x\)。

	将其代入度量拉回式，得到
	\[
		\mathrm{d}x^2+e^{2x^2}(Y_x\mathrm{d}x+Y_y\mathrm{d}y)^2
		=\mathrm{d}x^2+e^{2x^2}\mathrm{d}y^2.
	\]
	比较 \(\mathrm{d}x^2\) 系数，\(Y_x=0\)；比较 \(\mathrm{d}y^2\) 系数，\(Y_y^2=1\)。在一个小连通邻域内，连续且不为零的 \(Y_y\) 符号恒定。Jacobian 行列式恰为 \(Y_y\)，所以保持定向时有 \(Y_y=1\)，从而 \(\mathrm{d}(Y-y)=0\)，得到 \(Y=y+c\)。反转定向时有 \(Y_y=-1\)，同理得到 \(Y=-y+c\)。

	反向代入即可核验：这两类映射都保持 \(x\) 和 \(\mathrm{d}y^2\)，故保持度量，Jacobian 的符号分别为正、负。局部计算因而确定了全部局部自等距，而非只找到其中一族。
\end{solution}

\begin{exercise}\label{g1:ex:cylinder-local-isometry-covering}
	给定 \(\ell>0\) 与整数 \(m\ge2\)。对平圆柱
	\[
		C_{m\ell}=\mathbb R^2/(m\ell\mathbb Z\times\{0\}),\qquad
		C_\ell=\mathbb R^2/(\ell\mathbb Z\times\{0\}),
	\]
	考察映射 \(F([x,y]_{m\ell})=[x,y]_\ell\)。证明它良定义，是保持定向的局部等距及 \(m\) 叶覆盖，但不是整体等距同构。计算源圆柱上绕行一周的环路映到目标后的绕行次数，并写出两侧展开单值的对应关系。说明即使两边的切丛平行移动单值都平凡，这个覆盖仍不可能是单射。
\end{exercise}
\begin{solution}
	源中的两个代表元相差 \((km\ell,0)\)，它们在目标中也相差一个允许的平移，所以 \(F\) 良定义。两个商投影的局部逆将 \(F\) 表示为Euclid 平面中的恒等映射，故 \(F\) 光滑、保持定向且局部等距。

	对目标中的 \([x,y]_\ell\)，所有原像恰为
	\[
		[x+j\ell,y]_{m\ell},\qquad j=0,1,\ldots,m-1.
	\]
	它们两两不同，因为 \((j-k)\ell\) 不是非零的 \(m\ell\) 整数倍。取目标中由宽度小于 \(\ell\) 的小矩形给出的坐标域，它在源中的逆像分成这 \(m\) 个互不相交的小矩形，每一片由 \(F\) 同胚地映到原坐标域。因此 \(F\) 是 \(m\) 叶覆盖，且 \(m\ge2\) 说明它不单射。

	源中一周环路可写为 \(t\mapsto[x_0+tm\ell,y_0]_{m\ell}\)、\(0\le t\le1\)。在目标中它的平面提升仍从 \((x_0,y_0)\) 走到 \((x_0+m\ell,y_0)\)，故绕行次数为 \(m\)。相应地，基本群上的同态为 \(n\mapsto mn\)。若两边都取恒等的平面展开映射，则
	\[
		\rho_{m\ell}(n)(u,v)=(u+nm\ell,v)
		=\rho_\ell(mn)(u,v).
	\]
	两侧都有由 \(\partial_x,\partial_y\) 下降的整体平行标架，所以切丛平行移动单值平凡；展开单值却分别记录周期 \(m\ell\) 与 \(\ell\)。它们的最短不可缩环路长度不同，由\cref{g1:prop:cylinder-period-length}，甚至不存在其他整体等距把这两个圆柱识别。即使存在这样的局部等距与有限叶覆盖，两个圆柱的整体周期仍有差异。
\end{solution}
